\section{Algebraic reduction and stability of the torus pencil}
\label{sec:E}

\begin{statement}{E1}{Proved}{Vertical divisors and the field of meromorphic functions}
Let $Y$ be a connected compact complex manifold of dimension $n\geq2$
with $H^2(Y,\mathbb R)=0$. Let $g:Y\to C$ be a surjective holomorphic
map to a connected compact Riemann surface, with connected fibres.
Suppose that its smooth fibres over a nonempty open subset of $C$ are
K\"ahler. Every effective divisor of $Y$ is supported on finitely many
fibres, and pullback induces an isomorphism of fields
\[
 g^*: \mathcal M(C)\xrightarrow{\ \sim\ }\mathcal M(Y).
\]
In particular $a(Y)=1$. Reducedness of the other fibres is unnecessary.
\end{statement}
\begin{proof}
For a prime divisor $D$, the image $g(D)$ is analytic by properness. It
is a point or all of $C$. In the latter case choose a regular value $c$
in the stated open subset for which $F=g^{-1}(c)$ is not contained in
$D$. The canonical section of $\mathcal O_Y(D)$ restricts to a nonzero
section with nonempty effective zero divisor $D_F$ on $F$. For a
K\"ahler form $\omega$ on $F$, the divisor class formula gives
\[
 0=\int_F c_1(\mathcal O_Y(D)|_F)\wedge\omega^{n-2}
   =\int_{D_F}\omega^{n-2}>0.
\]
The first equality follows from the restriction homomorphism
$H^2(Y,\mathbb R)\to H^2(F,\mathbb R)$; the last follows from positivity
on every component with its positive multiplicity. Thus $g(D)$ is a
point. A divisor has finitely many prime components, proving the first
assertion.

For $h\in\mathcal M(Y)$ remove from $C$ the finitely many values
supporting its polar divisor. On the inverse image of the complement,
$h$ is holomorphic: a meromorphic function on a manifold having no
codimension-one pole extends across its codimension-two indeterminacy
set by the local normality of the structure sheaf. Its restriction to
each compact connected fibre is constant. The identity
$g_*\mathcal O_Y=\mathcal O_C$, supplied by Stein factorization and
the normality of $C$, therefore gives a holomorphic descended function
there. To check meromorphic extension at the removed values, take the
closure of the graph of $h$ in $Y\times\mathbb P^1$ and apply the proper
map $g\times1$. Its image is a one-dimensional analytic subset of
$C\times\mathbb P^1$, generically the graph of the descended function.
Its component dominating $C$ has generic degree one, and hence is the
graph of a meromorphic function on $C$. Pullback equals $h$ on a dense
open subset and therefore in the meromorphic function field. Pullback
is injective because $g$ is surjective. This proves the field
isomorphism and its asserted transcendence degree.
\end{proof}

\begin{statement}{E2}{Proved}{The non-holomorphic reduction theorem}
Let $Y$ be a compact connected complex threefold with $b_2(Y)=0$ and
let $g:Y\dashrightarrow C$ be a nonconstant meromorphic map to a
compact nonsingular curve which does not extend holomorphically.
For every holomorphic vector bundle $V$ on $Y$, there is a dense open
subset of each component of $\operatorname{Pic}(Y)$ on which
\[
 H^q(Y,V\otimes L)=0\quad(0\leq q\leq3).
\]
Consequently $\chi(Y,V\otimes L)=0$ for every $L\in\operatorname{Pic}(Y)$,
and $\int_Yc_3(Y)=0$.
\end{statement}
\begin{proof}
We give the graph proof, including the comparison needed at the
exceptional locus. The corresponding resolution comparison in the
printed proof of CDP20, Theorem~2.1, is justified separately in the
audit following E17. The graph method is that of \cite{LRS2019}, in the theorem labelled
\texttt{theorem main2 new} in the supplied Lehn--Rollenske--Schinko
source; the multiplicity filtration below specifies its use on a full
scheme-theoretic fibre.

The curve $C$ is rational. Indeed resolve the map by successive blowups
with smooth centres. A map from a projective space to a positive-genus
curve is constant, so a holomorphic map from the last blowup to such a
curve is constant on its exceptional projective-space fibres. Its graph
then descends to the preceding smooth space: a proper bimeromorphic
graph with singleton fibres over a normal space is an isomorphism.
Descending at each stage would extend $g$ on $Y$, contrary to its
definition. Take the Stein factor of a smooth resolution of the map
and replace $g$ by the induced meromorphic map to that factor. It
is still non-holomorphic, since otherwise the original map would
be holomorphic. The same descent argument makes its target
$\mathbb P^1$. Its general fibre is irreducible and reduced;
its image in the graph below need not be smooth.

Let $Z\subset Y\times\mathbb P^1$ be the closure of its graph and write
\[
 \pi:Z\to Y,\qquad h:Z\to\mathbb P^1
\]
for the holomorphic projections. Locally the map is $[a:b]$, where
$a,b$ are holomorphic germs with no common prime factor. The graph has
equation $av-bu=0$ in coordinates $[u:v]$ on $\mathbb P^1$; this
hypersurface is integral. On this local product the exact hypersurface
sequence is
\[
 0\to\mathcal O(-1)\xrightarrow{\ av-bu\ }\mathcal O
 \to\mathcal O_Z\to0,
\]
with the harmless line bundle from $Y$ included in the first term when
the local equations are changed. The cohomology of $\mathcal O$ and
$\mathcal O(-1)$ on $\mathbb P^1$ gives
$R\pi_*\mathcal O_Z=\mathcal O_Y$. Moreover $R^q\pi_*\mathcal F=0$
for $q\geq2$ and every coherent $\mathcal F$, since every fibre of
$\pi$ has dimension at most one. For every closed subspace $T\subset Z$,
the ideal sequence and these two facts imply
$R^1\pi_*\mathcal O_T=0$. These are statements about this graph
projection, with their explicitly given local equations.

The two sections defining the meromorphic map are sections of a line
bundle $P$ on $Y$ with no common divisorial component. Their pullbacks
generate an invertible ideal on $Z$. It defines an effective Cartier
divisor $\Sigma$ with
\begin{equation}\label{eq:E-graph-pencil}
 h^*\mathcal O_{\mathbb P^1}(1)=\pi^*P\otimes\mathcal O_Z(-\Sigma).
\end{equation}
The local ideal is generated by $a/u$ or $b/v$ on the corresponding
graph chart, proving both the asserted invertibility and the displayed
comparison. Every component of $\Sigma$ dominates $\mathbb P^1$:
its support lies over the codimension-two common-zero locus of $a,b$,
and the fibre over each such point is the entire projective line.
In particular $\Sigma$ meets every fibre. A fibre $Z_c$ is identified
by $\pi$ with the Cartier divisor $(a v_c-bu_c=0)$ in $Y$, including
its multiplicities. It is therefore Cohen--Macaulay and Gorenstein.

We recall the elementary vanishing used below. On a compact reduced
irreducible space $D$, for a locally free sheaf $W$ and a nonzero
effective Cartier divisor $A$, the descending finite-dimensional
spaces $H^0(D,W(-kA))$ stabilize. A section in their stable value
vanishes to every order along $A$, hence vanishes by the Krull
intersection theorem and the identity theorem. The stable value is
zero. On a Cartier divisor with multiplicities, successively remove
one prime component at a time. The corresponding exact sequences
have graded factors $\mathcal O_D(-D_{\rm previous})$ on reduced prime
components; tensoring by an invertible sheaf and a vector bundle
preserves exactness. Thus it suffices to apply the preceding vanishing
to the finitely many reduced-component factors.

We now choose $M\in\operatorname{Pic}(Y)$ and $k>0$ such that
\begin{equation}\label{eq:E-all-fibres}
 H^0\bigl(Z_c,\pi^*(V^*\otimes M^{-1})\otimes
       \omega_{Z_c}\otimes\mathcal O_{Z_c}(-k\Sigma)\bigr)=0
 \quad\text{for every }c.
\end{equation}
Here is a simultaneous construction. There are only finitely many
reducible or nonreduced fibres. Form the graph of their reduced prime
components, joining two which intersect in $Y$. Give a component
weight equal to its shortest distance from a component meeting
$\Sigma$ on $Z$. Every weight is finite because the fibres are
connected and meet $\Sigma$. Let $D_i$ be the sum, with coefficient
one, of the components of weight $i$. The inclusion $D_i\hookrightarrow
Y$ is a Cartier divisor since $Y$ is smooth.

Put $M=\mathcal O_Y(\sum_{i<d}n_iD_i)$, where $d$ is the largest
weight. Choose $n_{d-1}$ first, then $n_{d-2}$, and continue down to
$n_0$. On each component of weight $i>0$, $D_{i-1}$ restricts to a
nonzero effective Cartier divisor; $D_j$ with $j<i-1$ is disjoint
from it. After the higher-index coefficients have been fixed, the
preceding vanishing permits a choice of $n_{i-1}$ which annihilates
the sections on every multiplicity-filtration factor belonging to
weight $i$. Later choices do not affect those factors. The remaining
weight-zero factors are annihilated by increasing $k$, because
$\Sigma$ restricts to a nonzero effective Cartier divisor on each.
This proves \eqref{eq:E-all-fibres} on the finitely many exceptional
fibres.

For clarity, uniformity on the other fibres does not presume local
freeness of any higher direct image. First choose a power of
$-\Sigma$ annihilating sections on one general irreducible reduced fibre. Serre
duality and top-degree base change show that
$R^2h_*(\pi^*(V\otimes P^k))$ then has finite support. Increasing $k$
can only decrease that support, because multiplication by the
nonzero divisor cutting out $\Sigma$ injects the dual section spaces.
At each of its finitely many irreducible fibres the same vanishing
lemma, with its multiplicity filtration, kills the remaining group.
The effective twist $M$ cannot introduce a section on another fibre:
each of its prime divisors restricts there to an effective Cartier
divisor. This completes the simultaneous construction.

Restriction of \eqref{eq:E-graph-pencil} to $Z_c$ gives the actual line
isomorphism $\pi^*P|_{Z_c}=\mathcal O_{Z_c}(\Sigma)$. Serre duality
on the Gorenstein surface $Z_c$ turns \eqref{eq:E-all-fibres} into
\[
 H^2(Z_c,\pi^*(V\otimes M\otimes P^k)|_{Z_c})=0.
\]
Top-degree proper base change and Nakayama's lemma now give
$R^2h_*\pi^*(V\otimes M\otimes P^k)=0$. Serre vanishing on
$\mathbb P^1$ supplies $t$ for which
\[
 H^2\bigl(Z,\pi^*(V\otimes M\otimes P^k)
                       \otimes h^*\mathcal O(t)\bigr)=0.
\]
By \eqref{eq:E-graph-pencil}, this is
$H^2(Z,\pi^*(V\otimes M\otimes P^{k+t})(-t\Sigma))=0$.
The sheaf on $t\Sigma$ has zero second cohomology: its pushforward to
$Y$ is supported on the at-most-one-dimensional exceptional image,
and its higher direct images vanish by the graph calculation above
and the projection formula. The divisor exact sequence therefore
gives
$H^2(Z,\pi^*(V\otimes M\otimes P^{k+t}))=0$. Finally
$R\pi_*\mathcal O_Z=\mathcal O_Y$ and the projection formula give
an isomorphism with $H^2(Y,V\otimes M\otimes P^{k+t})$.

Since $H^2(Y,\mathbb Z)$ is finite, all exponents in this construction
can be chosen divisible by its exponent. The resulting twisting line
lies in $\operatorname{Pic}^0(Y)$. The same construction applied to
$V\otimes L_*$ reaches any fixed component $L_*\operatorname{Pic}^0(Y)$.
For degree zero use the vanishing lemma directly with any nonzero
effective fibre divisor on $Y$ and large negative multiples.
For degrees three and one use Serre duality, respectively with these
degree-zero and degree-two arguments applied to $V^*\otimes K_Y$.
Semicontinuity in the holomorphic exponential family of line bundles
over $H^1(Y,\mathcal O_Y)$ gives proper analytic jumping loci; their
finite union has dense open complement. The exponential family is
constructed by exponentiating a basis of additive Cech one-cocycles,
so it includes every line in $\operatorname{Pic}^0(Y)$ and the
semicontinuity assertion needs no compactness of that group.

All line bundles have rationally zero first Chern class. Hence
Riemann--Roch makes $\chi(V\otimes L)$ independent of $L$, and the
just-proved vanishing makes it zero. For $V=T_Y$, the terms involving
$c_1(Y)$ vanish in the index, leaving
$\chi(T_Y)=\tfrac12\int_Yc_3(Y)$. This proves the last assertion.
\end{proof}

\begin{statement}{E3}{Proved}{Excluded algebraic dimensions}
A positive-dimensional compact Moishezon manifold has $b_2>0$.
A compact complex manifold homeomorphic to $S^6$ has algebraic
dimension $0$ or $1$; in the latter case its algebraic reduction is
holomorphic onto $\mathbb P^1$.
\end{statement}
\begin{proof}
For the first assertion take a projective modification
$\pi:\widehat Y\to Y$. Choose a general very ample divisor
$\widehat D$ and a general complete-intersection curve
$\widehat C$ in the smooth projective manifold $\widehat Y$ so that
they meet outside the exceptional locus and their intersection is
finite. Their images $D,C$ meet properly and nontrivially. Since $Y$
is smooth, $D$ is Cartier and its effective restriction to the
normalization of $C$ has positive degree. Thus
$\langle c_1(\mathcal O_Y(D)),[C]\rangle>0$, proving $b_2(Y)>0$.
This is the proof of CDP98, Lemma~1.4 in the attached arXiv source;
the prompt's reference to Corollary~1.2 does not match that source's
numbering.

On a six-sphere the Euler number is $2$. Statement E2 consequently
excludes every meromorphic function with an indeterminate map to
$\mathbb P^1$. If two meromorphic functions were algebraically
independent, both would extend to holomorphic maps to
$\mathbb P^1$. Their product map to $\mathbb P^1\times\mathbb P^1$
has rank two at some point. In local target coordinates $u,v$
centred there, the meromorphic function $u/v$ pulled back to $Y$
has an indeterminacy: its restrictions to the local curves
$u=0$ and $v=0$ have distinct limits. This contradicts E2.
Thus $a(Y)\leq1$. For $a(Y)=1$, the algebraic reduction cannot be
non-holomorphic by E2, and its base has genus zero because pullback
injects $H^1(C,\mathbb Q)$ into $H^1(Y,\mathbb Q)=0$.
\end{proof}

\begin{statement}{E4}{Proved}{Degree-one cohomological flatness over a curve}
For a proper surjective holomorphic map $g:Y\to C$ from a connected complex
manifold to a nonsingular curve with connected fibres,
$R^1g_*\mathcal O_Y$ is locally free and
$H^0(Y_c,\mathcal O_{Y_c})=\mathbb C$ for every scheme-theoretic fibre.
\end{statement}
\begin{proof}
We rewrite \cite[Theorem labelled \texttt{raynaud}]{Viaclovsky2026Combs}.
Work over a disc with coordinate
$t$. Multiplication by $t$ is injective on $\mathcal O_Y$, so the map
is flat. For a proper map $h:Z\to\Delta_u$ with connected reduced
central fibre, $h_*\mathcal O_Z=\mathcal O_{\Delta_u}$, and $u$ a
nonzero divisor, the exact sequence for multiplication by $u$ shows
that $R^1h_*\mathcal O_Z$ has no $u$-torsion: the preceding map
$\mathcal O_{\Delta_u,0}\to H^0(Z_0,\mathcal O_{Z_0})=\mathbb C$
is surjective.

Resolve the divisor $g^{-1}(0)$ by a proper modification
$\mu:\widetilde Y\to Y$ with nonsingular source, so its pullback has
simple normal-crossing support and local equation
$t=z_1^{m_1}\cdots z_r^{m_r}$. Choose $M$ divisible by every
multiplicity, make $t=u^M$, and normalize. Denote the resulting
normal space by $Z$ and its map to $\Delta_u$ by $h$. Its central
fibre is reduced. Indeed, putting $e_i=M/m_i$, each local normal
component is the quotient of a polydisc by
\[
 G_0=\{(\zeta_i):\zeta_i^{e_i}=1,\ \prod\zeta_i=1\},
 \quad z_i=x_i^{e_i},\quad u=\zeta\prod x_i,
\]
where the appropriate constant root of unity $\zeta$ specifies the
chosen component. This includes the distinct components when the
exponents have a common divisor. This finite normal model is
birational onto that component.
For its power-series ring $R$, averaging gives
$R^{G_0}/uR^{G_0}=(R/uR)^{G_0}$, which is reduced. The Stein factor
of $h$ has generic degree one, since the general fibre is connected;
normality makes it $\Delta_u$. Thus $Z_0$ is connected and the
preceding argument gives local freeness of $R^1h_*\mathcal O_Z$.

The deck group $G=\mu_M$ has quotient $Z/G=\widetilde Y$.
To verify this, the invariant subring of the integral normalization
of a local base-change algebra lies in the invariant fraction field
and is integral over the normal ring of $\widetilde Y$; hence it is
that ring. Exactness of finite pushforward and of averaging gives
\[
 R^1(g\mu)_*\mathcal O_{\widetilde Y}
   =\bigl(b_*R^1h_*\mathcal O_Z\bigr)^G,
 \qquad b(u)=u^M.
\]
The right side is a direct summand of a locally free sheaf, since
$b$ is finite flat. The degree-one Leray sequence for $g\mu$ injects
$R^1g_*\mathcal O_Y$ into this sheaf, because
$\mu_*\mathcal O_{\widetilde Y}=\mathcal O_Y$. It is torsion-free,
and therefore locally free on the nonsingular curve. Finally pushing
forward $0\to\mathcal O_Y\xrightarrow{t}\mathcal O_Y\to
\mathcal O_{Y_0}\to0$ and using the absence of $t$-torsion proves
$H^0(Y_0,\mathcal O_{Y_0})=\mathbb C$.
\end{proof}

\begin{statement}{E5}{Proved}{The infinitesimal map and the rank-two boundary}
Let $g:Y\to C$ be a proper connected-fibre map from a compact complex
threefold whose smooth general fibre is a two-dimensional complex
torus. There is a canonical exact sequence of vector spaces
\begin{equation}\label{eq:E-Leray-map}
0\to H^1(C,T_C)\xrightarrow{g^*}H^1(Y,g^*T_C)
 \to H^0(C,T_C\otimes R^1g_*\mathcal O_Y)\to0.
\end{equation}
On the smooth locus $U\subset C$, the relative tangent sequence
induces a sheaf morphism
\[
 \delta:T_C\otimes R^1g_*\mathcal O_Y|_U
                  \longrightarrow R^2g_*T_{Y/C}|_U.
\]
At a smooth torus fibre $F$, write its Kodaira--Spencer tensor as
$\kappa=u_1\otimes\eta_1+u_2\otimes\eta_2$, where
$\eta_1,\eta_2$ is a basis of $H^1(F,\mathcal O_F)$ and
$u_i\in\operatorname{Lie}(F)$. Then
\[
 \delta(a\eta_1+b\eta_2)
   =(bu_1-au_2)\otimes(\eta_1\wedge\eta_2).
\]
Its rank is the rank of $\kappa$. When this rank is two at one
regular value, the image of $H^1(Y,T_Y)\to H^1(Y,g^*T_C)$ lies
in $g^*H^1(C,T_C)$. For $C=\mathbb P^1$ this image is zero.
\end{statement}
\begin{proof}
The projection formula gives
$R^qg_*g^*T_C=T_C\otimes R^qg_*\mathcal O_Y$, and the coherent
Leray spectral sequence has only columns zero and one. Its
degree-one edge sequence is exactly \eqref{eq:E-Leray-map}.
On $U$ the sequence
$0\to T_{Y/C}\to T_Y\to g^*T_C\to0$ is exact. Its extension
class on $F$ is the Kodaira--Spencer tensor. The connecting map is
cup product with that class; applying alternating multiplication to
the two displayed basis vectors gives the formula and its rank.
Exactness of the higher-direct-image sequence says that the image
of a global tangent class restricts to the kernel of $\delta$.
Rank two holds on an open dense subset after it holds at one point.
By E4 the source $T_C\otimes R^1g_*\mathcal O_Y$ is locally free,
so a section zero on that subset is zero everywhere. This proves
the image assertion and the final statement.

At a critical point the sheaf map $dg$ need not be onto. Its correct
image is $\mathcal I_{\rm crit}\otimes g^*T_C$, giving
\[
0\to T_{Y/C}\to T_Y\to
 \mathcal I_{\rm crit}\otimes g^*T_C\to0.
\]
For $t=z^m$, the image ideal is $(z^{m-1})$; for $t=xy$ it is
$(x,y)$; for $t=xyz$ it is $(xy,xz,yz)$. On a root cover of a
smooth multiple fibre the ordinary tangent sequence for $z$ applies.
At a normal-crossing fibre the logarithmic map sends
$\sum a_i z_i\partial_{z_i}$ to
$(\sum a_i)t\partial_t$ and is onto the logarithmic base tangent
line. These formulas type the comparison at the special fibres.
The proof above only uses the genuine smooth-locus boundary and
the globally locally free source, so it makes no assumption about
local freeness of $R^2g_*T_{Y/C}$ at a special value.
\end{proof}

\begin{statement}{E6}{Proved}{Stability of a rank-two torus fibration}
Let $g:Y\to\mathbb P^1$ be as in E5 and suppose its
Kodaira--Spencer tensor has rank two at one regular value. For every
proper deformation $\mathcal Y\to(S,0)$ of $Y$ over a complex
analytic germ, including a nonreduced germ, there is, after shrinking
$S$, a holomorphic map
\[
 G:\mathcal Y\longrightarrow\mathbb P^1\times S
\]
over $S$ whose central restriction is $g$. Every $G_s$ is surjective
with connected fibres and smooth general fibre a complex two-torus.
With $b_2(Y)=0$, every sufficiently small deformation has algebraic
dimension one.
\end{statement}
\begin{proof}
We prove the lifting assertion at every Artin order; this is more
than the tangent-space assertion of E5. The obstruction calculations
below use the Cech construction of holomorphic-map deformations
in \cite[\S4, equations (4.3)--(4.5), and \S6]{Horikawa1974};
we spell out the portions needed here. The injectivity hypothesis
on degree-two tangent cohomology in its Theorem~6.1 is replaced
here by the explicitly proved unobstructedness of the torus source.

First fix a rank-two regular fibre and its embedding $i:F\hookrightarrow
Y$. The normal line is canonically $N=\mathcal O_F\otimes T_c\mathbb
P^1$. The tangent sequence and E5 imply that
\begin{equation}\label{eq:E-surj-embedding}
 H^1(F,T_F)\longrightarrow H^1(F,i^*T_Y)
 \quad\hbox{is surjective}.
\end{equation}
Let $A'\to A$ be a small extension of local Artin complex algebras,
with kernel $I$ annihilated by the maximal ideal. Suppose an embedding
$i_A:F_A\hookrightarrow Y_A$ deforming $i$ and a deformation
$Y_{A'}$ lifting $Y_A$ have been specified. The deformation functor
of a complex torus is smooth: with a real lattice marking its small
complex structures are graphs of arbitrary small complex linear maps
between the two complementary Hodge subspaces; constant Beltrami
tensors have zero bracket. This gives a smooth four-dimensional
miniversal germ for $F$. In particular $F_A$ has a lift $F_{A'}$.

Choose local lifts of $i_A$ on sufficiently small coordinate charts.
They exist because the target deformation is smooth over $A'$.
Their differences on overlaps form a one-cocycle in
$i^*T_Y\otimes I$. Changing the lifted transition functions of
$F_{A'}$ by a cocycle $v_{ij}\in T_F\otimes I$ changes this
obstruction by $di(v_{ij})$. Such changes are precisely the action
of $H^1(F,T_F)\otimes I$ on the set of lifts of $F_A$; their
cocycle law has no quadratic term because $I\mathfrak m_{A'}=0$.
Surjectivity \eqref{eq:E-surj-embedding} therefore makes the map
obstruction zero. A zero-coboundary correction glues the local maps.
The resulting proper map is a closed immersion: the map on local
rings is surjective modulo the nilpotent ideal, hence surjective
by Nakayama, and its underlying central map is the embedding $i$.
Thus the given embedded torus lifts in every ambient small extension.

Now suppose $g_A:Y_A\to\mathbb P^1_A$ has been lifted and that
$Y_{A'}$ is prescribed. Choose local lifts $t_i$ of the target
coordinate of $g_A$; on overlaps their discrepancies define the
holomorphic-map obstruction
\[
 \alpha\in H^1(Y,g^*T_{\mathbb P^1})\otimes I.
\]
Coordinate changes on the target give exactly the transition factors
of $g^*T_{\mathbb P^1}$, and the triple-overlap relation gives its
Cech cocycle law. A coboundary makes the local maps agree, so
$\alpha=0$ is sufficient as well as necessary.

For a rank-two regular value $c$, the divisor
$F_A=g_A^{-1}(c)$ is a smooth deformation of the central torus.
The preceding embedding argument provides a Cartier divisor
$F_{A'}\subset Y_{A'}$ lifting it. This forces the restriction
of $\alpha$ to $F$ to vanish, as can be checked without a comparison
of obstruction theories: choose local equations $t_i-\eta_i$ for
$F_{A'}$, where $\eta_i$ lies in $I\otimes\mathcal O_Y$.
Two such equations differ by a unit. Modulo the central equation
$t=0$, that unit difference disappears, and their overlap equality
reads
\[
 (t_i-t_j)|_F=\eta_i|_F-\eta_j|_F.
\]
This is exactly the assertion that the restricted obstruction
cocycle is a coboundary in $N\otimes I$.

By \eqref{eq:E-Leray-map} and $H^1(\mathbb P^1,T_{\mathbb P^1})=0$,
$\alpha$ is a section of the locally free sheaf
$T_{\mathbb P^1}\otimes R^1g_*\mathcal O_Y$, tensored with $I$.
Its evaluation at every rank-two regular fibre is zero. That set
is open dense; hence the section and $\alpha$ are zero. The map
therefore lifts across every small extension.

We pass from this formal statement to a holomorphic map over the
given analytic germ. The germ of maps from $\mathcal Y/S$ to
$\mathbb P^1\times S/S$ is the open graph locus of the relative
Douady space; graphs are closed analytic subspaces, proper and
flat over the parameter, and their projection to $\mathcal Y$ is
an isomorphism on this open locus. Douady representability
\cite[\S9, Theorems 1--2; Appendix, \S2, Theorem 1]{Douady1966} therefore
gives a finite-dimensional analytic germ $\mathcal H\to S$
through the point $[g]$. The just-proved lifting property for every
small Artin extension says that this morphism is formally smooth
at $[g]$. The local analytic smoothness criterion identifies a
formally smooth morphism of finite-dimensional analytic germs with
a projection $S\times(\mathbb C^r,0)\to S$ after shrinking. One
can see the criterion on presentations of their convergent local
algebras: lifting across all square-zero quotients makes the
relative conormal obstruction zero, and the analytic implicit
function theorem eliminates the corresponding equations. The
zero section of this product supplies the required holomorphic
family of graphs. This argument applies to nonreduced $S$.

At one regular value, the family $G$ remains a proper submersion
near the central compact fibre. Its nearby fibres are connected
deformations of a torus. Each map $G_s$ is nonconstant, and hence
surjective. Its Stein factorization has finite degree one because
it has a connected regular fibre in this neighbourhood. Thus
every fibre is connected. Over the connected smooth locus of
$G_s$, Ehresmann's theorem identifies the differentiable fibres
with $T^4$. Every compact complex surface with this differentiable
type is a torus: its even first Betti number makes it K\"ahler
by \cite[Theorem 11]{Buchdahl1999};
its Albanese has dimension two; the exterior algebra on its four
integral degree-one classes shows the Albanese map has degree one;
and a nontrivial exceptional curve of this degree-one map would
include a $(-1)$-curve, impossible for the even intersection form
of $T^4$. Hence the general fibre is a complex two-torus.
Finally $b_2$ is constant in the differentiable deformation family,
and E1 proves $a(Y_s)=1$.
\end{proof}

The relevant earlier comparison is \cite[Remark 4.1, Proposition 4.2,
Theorem 4.4]{Catanese2004}. Its cup-product description is the same
boundary used in E5. Its smoothness criterion assumes
$H^0(C,T_C\otimes R^1g_*\mathcal O_Y)=0$, and its principal-bundle
application allows the positive-genus base to deform. In the present
examples that space of sections is nonzero; the embedded-fibre
argument above proves that the actual obstruction classes in it
vanish. Searches on 15 November 2026 for ``torus fibration deformation
Kodaira--Spencer rank'', ``rank two torus fibration deformation'',
and the cited 2026 arXiv identifiers did not locate an earlier theorem
with the statement of E6. We make no priority claim.

\begin{statement}{E7}{Proved}{Application to the named six-sphere families}
Every member of the torus-fibred families listed in the prompt,
including the period family $X_c$, its members with $|p|>1$,
Viaclovsky's $X_{3k,4l}$, his two-parameter family and his combs,
and the Liu--Rollenske family, satisfies the rank hypothesis of E6.
Consequently every sufficiently small deformation retains a
holomorphic torus fibration over $\mathbb P^1$. Every such member
with $b_2=0$, and every sufficiently small deformation of it, has
algebraic dimension one. In particular the fibration-retention
question in Section~10(1) of the supplied comb paper has an
affirmative answer.
\end{statement}
\begin{proof}
We verify a local condition common to these actual constructions.
At their hexagonal fibre, a marking on the punctured disc puts
the Hodge periods in the form
\begin{equation}\label{eq:E-log-period}
  (I_2,Z(q)),\qquad
 Z(q)=\frac{\log q}{2\pi i}B+H(q),
 \qquad \det B\ne0,
\end{equation}
where $H$ is holomorphic at zero. Therefore
\[
 qZ'(q)=\frac{B}{2\pi i}+qH'(q)
\]
is invertible for all sufficiently small nonzero $q$. The
Kodaira--Spencer tensor of a marked torus with periods $(I,Z)$
is represented by $Z'$ after the invertible identifications of
the Lie algebra and $H^1(\mathcal O)$ with the chosen period
coordinates. Thus it has rank two at those regular fibres.

Here the existence and nonsingularity of $B$ are checked directly
in the supplied local models, including
\cite[Proposition labelled \texttt{hexfill}]{Viaclovsky2026Combs},
\cite[Proposition labelled \texttt{mum}]{Viaclovsky2026Two}, and
\cite[Proposition labelled \texttt{prop: compactification X}]{LiuRollenske2026}.
The AC cusp has logarithmic matrix
$B_0=\left(\begin{smallmatrix}0&1\\-1&0\end{smallmatrix}\right)$.
The affine twists and finite logarithmic transformations change
the finite collars and constant gluing data, while the punctured
cusp matrix still has this same invertible logarithmic coefficient.
Engel's two deck transformations at the $I_1$-fibre have independent
valuation vectors $(1,0)$ and $(0,1)$, again giving an invertible
$B$. In the two-parameter family the designated fibre over $p_3$
is the full-lattice rank-two filling; its displayed two valuation
vectors are independent. Liu--Rollenske's compactification over
$I_1$ uses the same two independent valuation directions. Each
comparison is a change of the integral marking and of the
holomorphic fibre frame, so preserves the rank of the
Kodaira--Spencer tensor.

For the combs, \cite[Proposition labelled \texttt{periodextension},
case (i)]{Viaclovsky2026Combs} writes a flat frame as
$e_i=\alpha_i+\ell(q)\beta_i$, $f_i=\beta_i$, with
$\ell(q)=\log q/(2\pi i)$, and the extended Hodge subbundle as
the rows of $(I,H(q))$ in this frame. The fixed value of the
smoothed period section at $q=0$ keeps it transverse to the
$f_i$-plane. Its rows therefore remain $(I,H_s(q))$ with
$H_s$ holomorphic. In the flat frame they are
$(I,\ell(q)I+H_s(q))$, which is \eqref{eq:E-log-period} with
$B=I$. This verifies the needed condition for every comb member,
without a formula for its global splitting type.

All the sources construct the indicated members as compact
connected-fibre holomorphic maps with smooth general two-torus
fibre. E6 applies, and E1 gives the algebraic-dimension assertion.
On the differentiable six-sphere these neighbourhoods of the
listed structures constitute an open set on which $a=1$.
\end{proof}

\begin{statement}{E8}{Proved}{The theta determinant and its exact rank locus for $X_c$}
Use the period convention
$\Omega=\left(\begin{smallmatrix}\tau&\mu\\\mu&\beta/6\end{smallmatrix}\right)$,
Jacobi's odd theta function with periods $1,\tau$, and
$\beta/6=(i/\pi)\log(\vartheta_1(\mu,\tau)/\eta(\tau))+\text{constant}$.
Writing $h=\partial_\mu\log\vartheta_1$, one has the identity
\begin{equation}\label{eq:E-theta-general}
 \det\frac{d\Omega}{d\tau}
 =-\left(\mu'-\frac{i}{2\pi}h\right)^2
       -\frac{\wp(\mu,\tau)}{4\pi^2}.
\end{equation}
For the actual AC period branch $\wp(\mu,\tau)=0$, this becomes
\[
 \det\frac{d\Omega}{d\tau}
       =-\frac{g_2(\tau)^2}{36\pi^2g_3(\tau)}.
\]
On the smooth root-family base, the Kodaira--Spencer tensor has rank
one at the order-three point and rank two at every other finite
point, including the order-four point. Its determinant has order
two at the order-three point. The logarithmic tensor at the cusp
has rank two.
\end{statement}
\begin{proof}
The symmetric matrix has a concrete comparison with the raw AC
periods. Adjoin $(\mu,\beta/6)$ to the lattice generated by
$(1,0),(0,1),(\tau,\mu),(6\mu,\beta)$. This enlarges the
lattice by index six; identity on $\mathbb C^2$ induces a finite
holomorphic isogeny of degree six to the torus with period matrix
$(I,\Omega)$. Its differential and pullback on complex first
cohomology are isomorphisms, so the Kodaira--Spencer ranks agree.
This is the comparison used in the determinant calculation.

The normally convergent theta series gives its heat equation
$\partial_\tau\vartheta_1=(4\pi i)^{-1}\partial_\mu^2\vartheta_1$.
The product expansion of $\eta$ gives
$\eta'/\eta=\pi i E_2/12$. Finally the elliptic function
$-\partial_\mu h-\pi^2E_2/3$ has the principal part and constant
term of $\wp$, and hence equals $\wp$. Substituting these three
identities into
$\det\Omega'=(\beta/6)'-\mu'^2$ gives
\eqref{eq:E-theta-general}. Thus the square identity in Lead~9
uses the equation $\wp(\mu)=0$; a general section requires the
displayed additional term unless its other period terms cancel it.

For the zero branch differentiate $\wp(\mu,\tau)=0$. Put
$c=\pi^2E_2/3$ so $h'=-\wp-c$. Differentiating the heat equation
once in $\mu$, using $\wp''=6\wp^2-g_2/2$, and using
$E_2'=(\pi i/6)(E_2^2-E_4)$, yields at $\wp=0$
\[
 \partial_\tau\wp=\frac{h\wp'}{2\pi i}
                       +\frac{i g_2}{6\pi}.
\]
Consequently
\[
 \mu'-\frac{i h}{2\pi}=-\frac{i g_2}{6\pi\wp'},
 \qquad (\wp')^2=-g_3,
\]
which proves the second formula wherever $g_3\ne0$. It is an
identity of meromorphic expressions and therefore extends across
the branch values in a local uniformizing parameter.

The equation $j(\tau)=1728t$ shows that $\tau-\rho$ has order one
at the order-three root parameter, while $\tau-i$ has order two
at the order-four root parameter. The modular form $g_2$ has a
simple zero at $\rho$ and $g_3$ a simple zero at $i$, with no
other zeros modulo the modular group. Multiplication by
$(d\tau/ds)^2$ therefore gives an order-two zero at the first
point, and cancels the order-two zero of $g_3(\tau(s))$ at the
second. The tensor is nonzero at the first point because
$d\tau/ds\ne0$ there, hence has rank exactly one. Away from the
two root points the coordinate change is unramified and the
formula has nonzero determinant. At the cusp
$\tau=(\log q)/(2\pi i)+O(q)$, $\mu$ has a finite limit, and
$\beta=-\tau+O(1)$; hence the logarithmic derivative has limiting
matrix $\operatorname{diag}(1/(2\pi i),-1/(12\pi i))$.
Its determinant is nonzero.
\end{proof}

\begin{statement}{E9}{Known}{The reduced one-dimensional Kuranishi germ of the original family}
For every admissible $c$, the original $(3,4)$ family satisfies
$h^\bullet(X_c,T_{X_c})=(1,1,1,0)$ and has a smooth reduced
one-dimensional Kuranishi germ. Its displayed period parameter is
a local miniversal parameter. These statements occur in
\cite{DGKX2026v1}, in the results labelled
\texttt{p2:thm:tangent-cohomology} and
\texttt{p2:thm:kuranishi-aut}; version two \cite{DGKX2026v2} omits them.
\end{statement}
\begin{proof}
We confirm the particular implication used here from the supplied
R2--R3 data. On $\mathfrak B=\mathbb P(3,4)$, both vertical and
structure-cohomology bundles are $E=G=\mathcal O\oplus\mathcal O(-1)$.
The vertical subline is generated by the invariant translation
field; its quotient has twelfth power $\mathcal O(-p_0)$ by the
discriminant's simple cusp zero, hence is $\mathcal O(-1)$ in the
torsion-free root Picard group. The extension splits because
$H^1(\mathcal O(1))=0$. For $G$, relative duality and the
normalization complex of the hexagon give
$\det G=R^2\bar f_*\mathcal O_X=\mathcal O(-1)$; the invariant
integral cocycle $\gamma$ is nonzero on every ordinary fibre and
on the dual-graph cycle at the cusp. It gives the same split
extension. This checks the nonvanishing at the cusp required in R3.

The horizontal logarithmic line is $\mathcal O(-5)$. Projection
and Leray applied to
$0\to\bar f^*E\to T_X(-\log W)\to\bar f^*\mathcal O(-5)\to0$
give logarithmic tangent cohomology $(1,1,0,0)$: every nontrivial
root line appearing is one of
$\mathcal O(-1),\mathcal O(-2),\mathcal O(-5),\mathcal O(-6)$,
whose coarse direct image is $\mathcal O_{\mathbb P^1}(-1)$,
while the two surviving trivial summands lie in degrees zero
and one. The branch-normal sequence has quotient
$\eta_*K_{\mathrm dP_6}$, with cohomology $(0,0,1)$; the local
map sends each ambient vector field to its separate normal
coefficient on each branch. Hence its exact cohomology sequence
gives $(1,1,1,0)$ for $T_X$.

The parameter derivative is nonzero. The fixed real marking of
the torus lattice gives the constant Beltrami class
\[
 -\frac{1}{2iD}\,\partial_{\zeta_2}\otimes
 (d\bar\zeta_2-a\,d\bar\zeta_1),
 \quad a=\frac{\operatorname{Im}\mu}{\operatorname{Im}\tau},
 \quad D=\operatorname{Im}\beta_c
       -\frac{6(\operatorname{Im}\mu)^2}{\operatorname{Im}\tau}.
\]
A nonzero constant Dolbeault tensor on a torus is not exact,
because integration of every derivative of a periodic coefficient
is zero. A first-order trivialization of $X_c$ would preserve
the relative complete pencil $|-2K|$. Its induced infinitesimal
automorphism of $\mathbb P^1$ fixes all three critical values:
the equations $t=s^3$, $t-1=s^4$, and $q=xyz$ make each critical
value a constant section over the dual numbers. An infinitesimal
projective transformation fixing three distinct points is zero.
The trivialization would therefore trivialize the nonzero fibre
class just displayed, a contradiction. Thus the period derivative
spans the one-dimensional $H^1(T_X)$.

Write the Kuranishi local algebra as $\mathbb C\{z\}/I$ in its
one-dimensional tangent space. The existing proper period family
defines a local classifying map $z\mapsto\phi(s)$ with
$\phi'(0)\ne0$. The analytic inverse-function theorem makes
$\mathbb C\{z\}\to\mathbb C\{s\}$ an isomorphism. Since it
annihilates $I$, necessarily $I=0$. This verifies smoothness and
reducedness as well as miniversality, and every sufficiently small
deformation is a member of that period family.
\end{proof}

\begin{statement}{E10}{Refuted}{A fixed positive-genus target in the unrestricted stability assertion}
The extension assertion of E4(ii) of the prompt, with an arbitrary
fixed curve $C$ in place of $\mathbb P^1$, is false even for a torus
fibration with Kodaira--Spencer rank two at a regular value.
\end{statement}
\begin{proof}
Let $S\to\mathbb P^1_t$ be the rational elliptic surface obtained
by blowing up the nine base points of the Hesse pencil
$x^3+y^3+z^3-3txyz=0$. Its critical values are
$\{1,\omega,\omega^2,\infty\}$, where $\omega^3=1$, $\omega\ne1$.
Pull another copy back by $\phi(t)=1/(t-4)$; its critical values
are $\{4,5,4+\omega,4+\omega^2\}$, disjoint from the first set.
Their fibre product $Z\to\mathbb P^1$ is a smooth projective
threefold, since at every point at least one elliptic-surface
map is a submersion. Its general fibre is a product of two
elliptic curves. Both elliptic $j$-functions are nonconstant,
so their two Kodaira--Spencer derivatives are simultaneously
nonzero off a finite set. The tensor of their product is then
diagonal of rank two.

Let $C_s$ be the smooth genus-two double cover
\[
 y^2=(t-10)(t-11)(t-12)(t-13)(t-14)(t-15-s)
\]
of $\mathbb P^1$, for $s$ in a small disc. All its branch values
avoid the critical values of $Z$. Thus
$Y_s=Z\times_{\mathbb P^1}C_s$ is a smooth projective deformation
of a compact threefold, and $g_s:Y_s\to C_s$ has connected torus
general fibre and rank two at a regular value. The Kodaira--Spencer
class $\theta_C\in H^1(C_0,T_{C_0})$ is nonzero: any infinitesimal
isomorphism of the hyperelliptic covers descends through their
unique hyperelliptic involutions to an infinitesimal projective
transformation fixing five distinct branch values; it is zero,
whereas the sixth branch value moves with derivative one.

All the cohomology used for the obstruction is explicit. Put
$L=\pi_C^*\mathcal O_{\mathbb P^1}(1)$ on $C_0$; the
double-cover canonical formula gives $K_{C_0}=L$, of degree two.
For either rational elliptic surface, relative duality gives
$R^0\pi_*\mathcal O=\mathcal O$ and
$R^1\pi_*\mathcal O=\mathcal O(-1)$: its relative canonical
line is $\pi^*\mathcal O(1)$ and every fibre has
$h^\bullet(\mathcal O)=(1,1)$. Relative Kunneth and the flat
base change to $C_0$ consequently give
\[
 R^qg_*\mathcal O_Y=
 \begin{cases}
  \mathcal O_{C_0}&q=0,\\
  L^{-1}\oplus L^{-1}&q=1,\\
  L^{-2}&q=2,\\
  0&q>2.
 \end{cases}
\]
The vertical tangent sheaf is $g^*(L^{-1}\oplus L^{-1})$:
the invariant vector on each elliptic factor generates its
relative tangent sheaf, also at a node, where the kernel of
$d(xy)$ is generated by $x\partial_x-y\partial_y$.
Riemann--Roch on the genus-two curve gives
$h^0(L^{-a})=0$ and $h^1(L^{-a})=2a+1$ for $a>0$.
Projection and Leray now compute all degrees:
\[
\begin{array}{c|rrrr}
 \mathcal F&h^0&h^1&h^2&h^3\\\hline
 \mathcal O_Y&1&2&6&5\\
 g^*T_{C_0}&0&3&10&7\\
 T_{Y/C_0}&0&6&20&14
\end{array}
\]
All higher groups vanish. In particular
$H^0(C_0,T_{C_0}\otimes R^1g_*\mathcal O_Y)=0$ and the
injection $H^1(C_0,T_{C_0})\to H^1(Y,g^*T_{C_0})$ is an
isomorphism between three-dimensional spaces.

Compatibility of the Cech transition equations for the family
$g_s$ gives
\[
 dg(\theta_Y)=g^*\theta_C
       \quad\text{in }H^1(Y_0,g^*T_{C_0}).
\]
The right side is nonzero because the first map of
\eqref{eq:E-Leray-map} is injective. Hence the first-order
obstruction to extending $g_0$ with fixed target $C_0$ is nonzero.
This proves the counterexample with an actual smooth proper
family and its obstruction class. The obstruction belongs to
the base-deformation summand specified exactly in E5.
\end{proof}

\begin{statement}{E11}{Proved}{All coprime $(3k,4l)$ multiplicities}
Fix admissible logarithmic-transformation data with multiplicities
$r=3k$, $s=4l$, where $k,l>0$ and $\gcd(r,s)=1$, in the family
constructed from the original elliptic-surface data. For every member
which is diffeomorphic to $S^6$, as in
\cite[Theorem labelled \texttt{t:main}]{Viaclovsky2026Log},
\[
 h^\bullet(X,T_X)=(1,1,1,0),
\]
its Kuranishi germ is a smooth reduced line, and varying the nonzero
linearization scale with those data fixed is locally miniversal.
For the original $(3,4)$ affine-twist family, the same statements
hold whenever $p\ne0$, including $|p|>1$.
\end{statement}
\begin{proof}
Put $\mathfrak B_{r,s}=\mathbb P(r,s)$. The root morphism
$\sigma:\mathfrak B_{r,s}\to\mathfrak B_{3,4}$ over the identical
coarse $\mathbb P^1$ sends the two root sections to their $k$th and
$l$th powers. In line notation
$\sigma^*\mathcal O_{\mathfrak B_{3,4}}(1)
=\mathcal O_{\mathfrak B_{r,s}}(kl)$.
The vertical transition matrices and the action on fibre
$\mathcal O$-cohomology depend only on the linear monodromy.
The extra affine translations act trivially on both. The same
cusp comparison used in E9 therefore gives the actual bundles
\[
 E=G=\mathcal O\oplus\mathcal O(-kl),\qquad
 R^2\bar f_*\mathcal O_X=\mathcal O(-kl),
\]
and
\[
0\to\bar f^*E\to\mathcal T:=T_X(-\log W)
 \to\bar f^*L\to0,
\quad L=\mathcal O(r+s-rs),
\quad K_X=\bar f^*\mathcal O(kl-r-s).
\]
For example the normal-cover coordinate at the first finite point
is $u=v^k$, and differentiation changes the ordinary root tangent
line from the old one to $T_{\mathfrak B_{r,s}}$; the vertical
bundle and its cohomology are pulled back through $u=v^k$.
This specifies the comparison rather than identifying the two
horizontal tangent lines. At the cusp $W$ remains the same
normal-crossing divisor, so its logarithmic base correction is
$-p_0=-rs$ and its normal quotient is $\eta_*K_{\mathrm dP_6}$.

Locally $\sigma_*\mathcal O=\mathcal O$ and its positive direct
images vanish: the additional cyclic group has invariants
$\mathbb C\{v\}^{\mu_k}=\mathbb C\{u\}$, and averaging is
exact. Consequently the cohomology of the pulled-back lines
$\mathcal O(-kl)$ and $\mathcal O(-2kl)$ is the cohomology of
the original $\mathcal O(-1)$ and $\mathcal O(-2)$, hence zero.
Projection and Leray give
$h^\bullet(X,\bar f^*E)=(1,1,0,0)$.
For the horizontal line, its coarse direct image is
$\mathcal O_{\mathbb P^1}(-1)$: a vector field on a root disc
must have its coefficient divisible by the root parameter,
so $\rho_*T_{\mathfrak B_{r,s}}=T_{\mathbb P^1}(-p_1-p_2)$,
and the extra $-p_0$ gives the asserted degree. Moreover
$L$ and $L\otimes\mathcal O(-kl)$ have negative orbifold degree,
so have no global sections. Thus
$H^0(X,\bar f^*L)=H^1(X,\bar f^*L)=0$.
The two exact tangent sequences now give
$h^0(X,T_X)=h^1(X,T_X)=1$.

We compute degree three independently. On the smooth torus locus,
the boundary on a relative holomorphic one-form is the transpose
of the Kodaira--Spencer tensor. The latter is generically invertible
by E7. Hence every ambient holomorphic one-form on an inverse-image
base disc is horizontal. It is the pullback of a base one-form.
At a finite root point the root map is a submersion, so that
one-form extends holomorphically in the root coordinate. At a
generic smooth point of $W$ its local equation is $q=z_1$, which
excludes a pole in its base coefficient. Therefore the canonical
pullback comparison gives
\[
 \bar f_*\Omega_X^1=\Omega_{\mathfrak B_{r,s}}^1.
\]
By Serre duality and the projection formula,
\[
 h^3(X,T_X)=h^0\bigl(\mathfrak B_{r,s},
                         \mathcal O(kl-2r-2s)\bigr)=0.
\]
For the last equality, a weighted monomial of that degree would
give nonnegative integers $a,b$ with
$3k(a+2)+4l(b+2)=kl$. Reduction modulo $k$ makes $b+2$ a positive
multiple of $k$, since $\gcd(k,4l)=1$. The second term alone
is then at least $4kl$, a contradiction. On a six-sphere
$\chi(T_X)=1$, giving $h^2=1$. For the original affine-twist
family, $r=3$, $s=4$, and the horizontal Leray terms are the
acyclic lines $\mathcal O(-5)$ and $\mathcal O(-6)$ on
$\mathbb P(3,4)$. Therefore $h^\bullet(\mathcal T)=(1,1,0,0)$,
and the quotient $\eta_*K_{\mathrm dP_6}$ gives
$h^\bullet(T_X)=(1,1,1,0)$ directly for every admissible affine
twist. This computation does not require the integral homology
to be that of a sphere.

The scale derivative changes the fourth torus period by a nonzero
constant multiple of the parameter derivative in E9, so its
restricted torus class is nonzero. A first-order trivialization of
the compact member would turn the given fibration into a
first-order deformation of the fixed map $f$ with fixed source.
Such deformations are $H^0(X,f^*T_{\mathbb P^1})
=H^0(\mathbb P^1,T_{\mathbb P^1})$, by connected fibres and
the projection formula. Thus they are precisely infinitesimal
projective transformations of the base. The two multiple critical
values and the hexagonal critical value remain the same three
constant sections in this scale family; the local equations
$t=v^r$, $t-1=v^s$, $q=xyz$ show this also over the dual numbers.
The base vector field is therefore zero. The nonzero torus class
cannot be trivialized, proving the nonzero global scale derivative.
The one-dimensional tangent space and the analytic local-algebra
argument at the end of E9 prove smoothness, reducedness and
miniversality.
\end{proof}

\begin{statement}{E12}{Refuted}{The universal local-freeness assertion of CDP98, Lemma 1.5}
Let $A$ be an elliptic curve with origin $0$, let
$P=\mathcal O_{A\times A}(\Delta-A\times\{0\}-\{0\}\times A)$,
and let $Y=A\times A\times\mathbb P^1$, with $g:Y\to A$
projection to the second factor. For
$V=P\boxtimes\mathcal O_{\mathbb P^1}(-2)$ one has
\[
 g_*V=R^1g_*V=0,\qquad R^2g_*V=\mathbb C_0,
 \qquad h^\bullet(Y,V)=(0,0,1,0).
\]
The sheaf in degree two has length one. For $t\ne0$ the fibre
cohomology is zero in every degree; for $t=0$ it is $(0,1,1)$.
\end{statement}
\begin{proof}
Each nontrivial degree-zero line on an elliptic curve has no
sections: a nonzero section would have an effective divisor of
degree zero and hence trivialize it. Riemann--Roch and Serre
duality then give zero first cohomology. The trivial line has
cohomology $(1,1)$. To compute the direct image with its length,
put $p:A\times A\to A$ and $M=P\otimes\operatorname{pr}_1^*
\mathcal O_A(0)$. Its fibres have degree one, so
$R^1p_*M=0$ and $p_*M$ is a line bundle. The exact sequence for
the section $\{0\}\times A$ gives the two-term complex
\[
 Rp_*P=\bigl[p_*M\longrightarrow\mathcal O_A\bigr].
\]
The arrow evaluates the unique section of the moving degree-one
line at the origin. Its zero is exactly $t=0$ and is simple:
in a local coordinate the moving zero divisor is the diagonal,
with equation $x-t$. Thus the arrow is the inclusion
$\mathcal O_A(-0)\to\mathcal O_A$, and
$Rp_*P=\mathbb C_0[-1]$. Tensoring with
$R\Gamma(\mathbb P^1,\mathcal O(-2))=\mathbb C[-1]$
gives $Rg_*V=\mathbb C_0[-2]$. This proves the direct-image,
length, and global-cohomology assertions. Fibre Kunneth gives
the stated fibre cohomology.

For any proper flat family over a nonsingular curve and a locally
free sheaf $V$, multiplication by a local parameter gives the
actual replacement exact sequence
\[
0\to (R^qg_*V)_c/t(R^qg_*V)_c
 \to H^q(Y_c,V|_{Y_c})
 \to (R^{q+1}g_*V)_c[t]\to0.
\]
It follows directly from the long direct-image sequence of
$0\to V\xrightarrow{t}V\to V|_{Y_c}\to0$. Here $[t]$ denotes
the kernel of multiplication by $t$. Local constancy of all
fibre cohomology dimensions implies local freeness and base
change by Grauert's theorem; generic vanishing alone permits
the torsion explicitly present above.
\end{proof}

In the supplied CDP98 arXiv source, Lemma~1.5 is used in
Proposition~2.5 to pass from a general-fibre vanishing to
$R^2f_*=0$, in the proof of Theorem~3.1 for the asserted
flatness of the two-parameter direct image, in Lemma~5.4 for
$R^1f_*\mathcal O$, in Theorem~5.6 to propagate sections from a
single fibre, and in Proposition~5.7 to remove a torsion
$R^2f_*\mathcal O$. Its first use is replaced in E2 by vanishing
on every full fibre; the special structure-sheaf assertion in
degree one is independently established in E4. The other uses
require their own hypotheses, rather than the refuted lemma.

\begin{statement}{E13}{Proved}{The Kuranishi germ of the two-parameter family}
For every member of the family of
\cite[Proposition labelled \texttt{family}]{Viaclovsky2026Two}, on the parameter
bidiscs with $\mu$ sufficiently close to zero and nonzero sufficiently
small linearization scale $\lambda$, one has
\[
 (h^0,h^1,h^2,h^3)(X_{\mu,\lambda},T_X)=(1,2,2,0).
\]
Its Kuranishi germ is a smooth reduced surface, and the displayed
$(\mu,\lambda)$-family is locally miniversal. In particular its
sufficiently small deformations all have algebraic dimension one.
\end{statement}
\begin{proof}
We work with the equation and section used in the construction,
\[
 S_\mu:\quad y^2=x^3+tx+\mu t+1,\qquad
 P_\mu=(-\mu,b),\qquad b^2=1-\mu^3,
\]
where the branch of $b$ is close to one. Write $p_1,p_2,p_3$ for its
three distinct nodal values, with the zero of the linearization at
$p_3$. The order-four good-reduction cover at infinity gives the
root curve $\mathfrak B=\mathbb P(1,4)$, with coarse projection
$\rho$ and $\rho^*\mathcal O_{\mathbb P^1}(1)=\mathcal O(4)$.
Let $\bar f:X\to\mathfrak B$ denote the root factorization and
$D=W_1+W_2+W_3$ the reduced divisor of Mumford fibres.

We spell out the bundles and the residue map, since their ranks alone
would not determine the answer. The vertical logarithmic Lie bundle
is
\[
 E=\mathcal O\oplus\mathcal O(-1).
\]
The first summand is generated by multiplication in the line-bundle
coordinate. For a concrete splitting of the second summand put
$\omega=dx/(2y)$ and use the relative differential
\begin{equation}\label{eq:E13-connection}
 \alpha=\frac{y+b}{x+\mu}\,\frac{dx}{2y}.
\end{equation}
Its residues are $1$ at $P_\mu$ and $-1$ at the zero section; there
is no pole at $-P_\mu$, because its numerator vanishes there.
Consequently, in the rational section of
$\mathcal O(P_\mu-O)$ having divisor $P_\mu-O$, the connection
$\nabla(hs)=(dh+\alpha h)s$ is holomorphic along each smooth
elliptic fibre. Here $s$ is that rational section and a point of
the line bundle is written $v s$; the horizontal lift of a vector
$V$ is therefore $\widetilde V=V-\alpha(V)v\partial_v$.
It is a multiplicative connection: the difference between its
pullback under multiplication and the sum of its two pullbacks is a
holomorphic one-form on the product of two elliptic curves, and its
restrictions to both coordinate axes vanish. Every holomorphic
one-form on that product is the sum of forms pulled back from those
axes, so this difference is zero. Its horizontal lift therefore
commutes with the lifted translation defining $X$.

At infinity put $t=r^{-4}$, $x=r^{-2}x'$, $y=r^{-3}y'$. Formula
\eqref{eq:E13-connection}, in the relative cotangent sheaf, becomes
\[
 \alpha=\frac{y'+br^3}{x'+\mu r^2}\frac{dx'}{2y'},
 \qquad y'^2=x'^3+x'+\mu r^2+r^6.
\]
This is the same holomorphic connection at the good-reduction
fibre; the limiting section $(0,0)$ is a nonzero two-torsion point,
so it remains disjoint from the zero section. Moreover
$\omega=r\,dx'/(2y')$, which gives the summand $\mathcal O(-1)$
for its dual Lie vector. At a nodal fibre the connection is
logarithmic at the two branches and its lift is a section of the
vertical logarithmic tangent bundle. Thus these calculations prove
the claimed splitting on the entire root curve, with the specified
comparison maps.

The inverse pairing associated with the nondegenerate alternating
class $\xi=u\wedge w+2\psi\wedge\delta$ in the construction induces a
holomorphic isomorphism
\[
 G:=R^1\bar f_*\mathcal O_X\simeq E.
\]
Here is the extension check. On a marked smooth fibre the periods
have cross terms in ratio $1:-2$, so the holomorphic one-form
subspace $F^1\subset H^1(F,\mathbb C)$ is Lagrangian for this pairing.
Pairing with $F^1$ therefore gives the isomorphism
$H^1(F,\mathbb C)/F^1\to(F^1)^*=\operatorname{Lie}(F)$.
At a Mumford point, subtract the logarithmic monodromy term from the
periods. In the resulting holomorphic frame the alternating pairing
is still constant and nondegenerate, and the two holomorphic
one-forms are the logarithmic coordinate differentials. Their dual
frame is the logarithmic Lie frame. The quotient frame gives
$R^1\bar f_*\mathcal O_X$: at a rank-one point its first class is
the class of the loop in the normalization diagram and its second
is the class on the elliptic normalization; at the rank-two point
both classes are loops in the normalization diagram. These are
bases, by the normalization exact sequence. Thus the pairing extends
as an isomorphism at these points. The finite good-reduction cover
gives the same conclusion at infinity, and descent uses invariance
of $\xi$. This identifies the comparison, including its value at
the special fibres. Cup product and the normalization sequences
give $R^2\bar f_*\mathcal O_X=\det G=\mathcal O(-1)$.

The logarithmic tangent and canonical sequences are therefore
\[
\begin{split}
 0&\longrightarrow\bar f^*E\longrightarrow T_X(-\log D)
 \longrightarrow\bar f^*\mathcal O(-7)\longrightarrow0,\\
 K_X&=\bar f^*\mathcal O(-4).
\end{split}
\]
Indeed $K_{\mathfrak B}=\mathcal O(-5)$ and
$T_{\mathfrak B}(-p_1-p_2-p_3)=\mathcal O(5-12)$; taking the
determinant of the displayed logarithmic sequence gives the
canonical formula after removing $D$ of degree $12$.
For root line bundles, cohomology is computed by weighted monomials
and Serre duality with $\mathcal O(-5)$. The projection formula and
Leray give
\[
\begin{split}
 h^\bullet(X,\bar f^*E)&=(1,1,0,0),\\
 h^\bullet(X,\bar f^*\mathcal O(-7))&=(0,1,2,1),\\
 h^\bullet(X,T_X(-\log D))&=(1,2,2,1).
\end{split}
\]
There are no uncomputed differentials in these equalities: the base
has dimension one, and the groups of the first displayed row in
degrees two and three vanish.

Let $Q=\bigoplus_{i=1}^3\eta_{i*}K_{\widetilde W_i}$, the normal
quotient of $T_X(-\log D)\to T_X$. Here
$\widetilde W_1,\widetilde W_2$ are ruled over elliptic curves and
$\widetilde W_3=dP_6$, so
\[
 h^\bullet(Q)=(0,2,3).
\]
It remains to compute the boundary
$H^1(Q)\to H^2(T_X(-\log D))$, between two-dimensional spaces.
Under Serre duality its dual is the residue homomorphism
\begin{equation}\label{eq:E13-residue}
 H^1\bigl(X,\Omega_X^1(\log D)\otimes K_X\bigr)
 \longrightarrow
 H^1(\widetilde W_1,\mathcal O)\oplus
 H^1(\widetilde W_2,\mathcal O).
\end{equation}
The dual logarithmic sequence, tensored with $K_X$, has horizontal
term $\bar f^*\mathcal O(3)$ and vertical quotient
$\bar f^*(\mathcal O(-4)\oplus\mathcal O(-3))$.
The latter has zero cohomology in degrees zero and one: the coarse
direct images of $\mathcal O(-4)$ and $\mathcal O(-3)$ are both
$\mathcal O_{\mathbb P^1}(-1)$, and every line in their tensor
products with $G$ has negative weighted degree, so has no section.
Hence the
source of \eqref{eq:E13-residue} is canonically
\[
 H^1(X,\bar f^*\mathcal O(3))
 =H^0\bigl(\mathfrak B,\mathcal O(3)\otimes G\bigr)
 =H^0(\mathcal O(3))\oplus H^0(\mathcal O(2))=\mathbb C^2.
\]
On the affine $t$-line the generators of these last two spaces are
constant in the frame just constructed. The residue in each
summand is a nonzero scalar times evaluation followed by
$G_{p_i}\to H^1(\widetilde W_i,\mathcal O)$.

Write $e_i$ for the double root in the nodal cubic at $p_i$.
The line killed by this last projection is the graph-loop class.
Under $G\simeq E$ it is the toric Lie line: in the local periods
$(q\alpha,\beta),(\gamma,\lambda\delta)$ it is the direction
of the degenerating first multiplicative coordinate, whereas the
second coordinate is the elliptic normalization. If $V$ is dual to
$\omega$, its horizontal lift $\widetilde V$ has vertical
line-coordinate component $-\alpha(V)$. At the nodal boundary put
\[
 a_i=\frac{b}{e_i+\mu}.
\]
Thus the toric line is generated by
$\widetilde V+a_i v\partial_v$, and the projection to the elliptic
Lie line sends $(v\partial_v,\widetilde V)$ to $(1,-a_i)$, up to
a common nonzero factor at that point. Consequently the matrix of
\eqref{eq:E13-residue}, after nonzero row rescalings, is
\[
 \begin{pmatrix}1&-a_1\\1&-a_2\end{pmatrix},\qquad
 a_1-a_2=\frac{b(e_2-e_1)}{(e_1+\mu)(e_2+\mu)}\ne0.
\]
The inequality is exact: $b\ne0$; distinct nodal values have
distinct double roots, since $p_i=-3e_i^2$; and $e_i=-\mu$ would
place $P_\mu$ at the node and force $b=0$. The boundary therefore
has rank two.

The rank-two cusp implies
$\bar f_*\Omega_X^1=\Omega_{\mathfrak B}^1$, by the generic
Kodaira--Spencer boundary and extension across the local models,
as in E11. It follows that
$h^3(T_X)=h^0(\Omega_X^1\otimes K_X)
=h^0(\mathcal O(-9))=0$. The long exact sequence with quotient
$Q$ now gives $h^0(T_X)=1$, $h^1(T_X)=2$, and $h^2(T_X)=2$.

Finally the constructed bidisc has injective Kodaira--Spencer map.
At $\mu=0$, writing $a_i=p_i(0)$, differentiation of
$4t^3+27(\mu t+1)^2=0$ gives $p_i'(0)=-9/(2a_i)$.
An infinitesimal automorphism fixing the multiple value at infinity
has the form $At+B$; the three equalities
$-9/(2a_i)=Aa_i+B$ would make a nonzero quadratic polynomial vanish
at all three $a_i$. The cross-ratio variation is therefore nonzero,
and remains so on a smaller disc. For fixed $\mu$, a normalized
double curve is $\mathbb C^*/\langle c(\mu)\lambda\rangle$.
Its elliptic modulus has nonzero $\lambda$-derivative when
$|c(\mu)\lambda|$ is sufficiently small, as follows from
$j(q)=q^{-1}+744+O(q)$. A trivial first-order deformation of $X$
would trivialize its intrinsic anticanonical pencil, its labelled
exceptional-value configuration, and this double curve; it can
therefore kill neither of the two independent tangent directions.
Thus the bidisc maps to the two-dimensional tangent space of the
Kuranishi germ with invertible differential. Present its local
ring as $\mathbb C\{z_1,z_2\}/I$; the induced substitution into
$\mathbb C\{\mu,\lambda-\lambda_0\}$ is a local analytic
isomorphism before quotienting, by the inverse function theorem.
Every element of $I$ vanishes under it, so $I=0$. This proves the
Kuranishi assertion. E1 and E6 give the algebraic-dimension assertion.
\end{proof}

\begin{statement}{E14}{Proved}{The Liu--Rollenske family and its transverse deformation}
Every member $Y(u)$ of the Liu--Rollenske family, with its scale
sufficiently small and nonzero as in its construction, satisfies
\[
 h^\bullet(Y(u),T_{Y(u)})=(1,2,2,0).
\]
Its Kuranishi germ is a smooth reduced surface. One parameter is
the given scale; a second parameter replaces its cuspidal fibre by
two rank-one nodal fibres. All sufficiently small deformations
retain a torus fibration and have algebraic dimension one.
\end{statement}
\begin{proof}
We derive the needed holomorphic invariants from the local models.
The commented-out holomorphic-invariants notes in the supplied
Liu--Rollenske source are not used as established results.

In the Weierstrass equation of \cite[Proposition labelled
\texttt{prop: surface facts}]{LiuRollenske2026}, make the explicit
coordinate comparison
\[
 t_{\rm new}=\frac{27}{4}(t_{\rm old}^{-1}-1),\quad
 x_{\rm new}=x_{\rm old}/t_{\rm old}^2,\quad
 y_{\rm new}=y_{\rm old}/t_{\rm old}^3.
\]
It gives $y^2=x^3+tx+t$, with section $P=(-1,i)$; the multiple
value is now infinity, the cusp is at zero and the remaining nodal
value is $-27/4$. Consider the explicit deformation
\begin{equation}\label{eq:E14-elliptic-family}
 y^2=x^3+tx+t+\epsilon,\qquad
 P_\epsilon=(-1,b),\quad b^2=\epsilon-1,
 \quad b(0)=i.
\end{equation}
Its finite discriminant is $4t^3+27(t+\epsilon)^2$.
The nodal value near $-27/4$ remains simple and holomorphically
labelled. The other two nodes coalesce to the type-II cusp at
$\epsilon=0$. Infinity keeps type $III^*$; on its order-four
good-reduction cover the equation is
\[
 y'^2=x'^3+x'+r^2+\epsilon r^6,
 \qquad P'_epsilon=(-r^2,br^3).
\]

These equations give a simultaneous construction of the compact
threefolds, including the central one. On a disc containing the cusp
and its two nearby nodal values, the total elliptic surface in
\eqref{eq:E14-elliptic-family} is smooth over $\epsilon$:
$\partial(y^2-x^3-tx-t-\epsilon)/\partial t=-(x+1)$ is a unit
near the cusp. Translation by $-2P_\epsilon$ is an automorphism
there. The line bundle $\mathcal O(P_\epsilon-O)$ and its lifted
translation vary holomorphically; its multiplier has a zero only
at the separate simple nodal value chosen for the hexagonal filling.
For a uniformly small nonzero scale its contraction on the
complement of the zero section is uniform on the compact elliptic
fibres. The quotient is consequently free, properly discontinuous,
proper over the disc, and smooth over the parameter. This quotient
itself fills the cusp and the nearby rank-one nodal fibres. At the
separate node the ordinary rank-two Mumford charts vary
holomorphically. At infinity the fixed order-four torsion translation
and good-reduction comparison vary holomorphically in the displayed
equation. Gluing along their marked punctured collars therefore
constructs a proper holomorphic family over $(\epsilon,u)$ whose
central scale slice is $Y(u)$. Every member has the central
differentiable type by Ehresmann. The supplied topology theorem
\cite[Theorem labelled \texttt{thm: main}]{LiuRollenske2026} thus
identifies these differentiable manifolds with $S^6$.

Write $\mathfrak B=\mathbb P(1,4)$ and retain its convention
$\rho^*\mathcal O(1)=\mathcal O(4)$. On the central threefold the
canonical form on the contracting quotient is the elliptic-surface
two-form wedged with $dv/v$. At the cusp it extends without a zero
or pole: from $t=(y^2-x^3)/(1+x)$ one gets
\[
 dt\wedge\frac{dx}{2y}=-\frac{dx\wedge dy}{1+x}.
\]
The finite-cover and hexagonal calculations are the same explicit
coordinate calculations as in E13. Since the elliptic surface has
$K_S=\pi^*\mathcal O_{\mathbb P^1}(-1)$, these comparisons give
\begin{equation}\label{eq:E14-canonical}
 K_Y=f^*\mathcal O_{\mathbb P^1}(-1)
     =\bar f^*\mathcal O_{\mathfrak B}(-4).
\end{equation}
The calculation is simultaneous in $\epsilon$. Adjunction on every
full fibre gives trivial dualizing sheaf. E4 gives $h^0=1$;
duality gives $h^2=1$; constancy of its Euler characteristic in the
proper flat family gives $h^1=2$. Thus, also on the two-dimensional
base $(t,\epsilon)$, Grauert's theorem gives locally free $R^qf_*\mathcal O$
with base change, for $q=0,1,2$.

The multiplicative connection
\[
 \alpha=\frac{y+b}{x+1}\frac{dx}{2y},\qquad
 \nabla(hs)=(dh+\alpha h)s
\]
has precisely the residues of the rational section $s$ with divisor
$P_\epsilon-O$. The same proof and good-reduction substitution as
in E13 give the logarithmic Lie bundle
$E=\mathcal O\oplus\mathcal O(-1)$, including at the central
cusp. There the vector dual to $dx/(2y)$ is the derivation
$2y\partial_x+3x^2\partial_y$ on the cuspidal fibre, and its
horizontal lift is well defined because
$(y+b)/(x+1)$ is regular near the singular locus.

We justify the identification with $G=R^1\bar f_*\mathcal O$
at the cusp without assuming a limiting-period rule there.
The inverse pairing associated with the invariant alternating
class gives $G\simeq E$ over the two-dimensional $(t,\epsilon)$
base away from its single cuspidal point: all remaining points
are smooth, nodal or the finite-cover points already computed in
E13. Both $G$ and $E$ are locally free across that point, by the
preceding cohomological-flatness and connection calculations.
The matrix entries of the isomorphism and its inverse extend
across that codimension-two set by Hartogs' theorem. Their products
are the identity on a dense open set, hence everywhere. Thus
\[
 G=\mathcal O\oplus\mathcal O(-1),\qquad
 R^2\bar f_*\mathcal O=\mathcal O(-1)
\]
on the central root curve as well; the second equality also follows
directly from relative duality and \eqref{eq:E14-canonical}.

Let $D=C+W$ be the sum of the cuspidal and hexagonal fibres.
For the local tangent calculation choose coordinates in which
$t=y^2-x^3$. They are obtained from the earlier formula by the
invertible change $\widetilde x=x/(1+x)^{1/3}$,
$\widetilde y=y/(1+x)^{1/2}$, with the branches equal to one at
zero, and then renaming the two new coordinates. The logarithmic
tangent module is freely generated by
\[
 2y\partial_x+3x^2\partial_y,\quad
 2x\partial_x+3y\partial_y,\quad\partial_z,
\]
where $z$ is a coordinate along the elliptic singular curve. The
second generator maps to $6t\partial_t$;
the first and third generate the relative logarithmic kernel.
Consequently the global exact sequence is
\[
 0\longrightarrow\bar f^*E\longrightarrow T_Y(-\log D)
 \longrightarrow\bar f^*\mathcal O(-3)\longrightarrow0,
\]
because $T_{\mathfrak B}(-C-W)=\mathcal O(5-8)$.
The vertical cohomology is $(1,1,0,0)$. Every line in the
horizontal Leray terms is $\mathcal O(-3)$ or $\mathcal O(-4)$,
which has no cohomology on $\mathbb P(1,4)$. Hence
\begin{equation}\label{eq:E14-log-coho}
 h^\bullet(T_Y(-\log D))=(1,1,0,0).
\end{equation}

The quotient at $W$ is $\eta_*K_{dP_6}$ with cohomology $(0,0,1)$.
At $C$, the map $v\mapsto v(t)|_C$ identifies the quotient with
the Jacobian ideal
\[
 J=(x^2,y)\subset\mathcal O_C.
\]
Indeed its kernel is precisely the logarithmic tangent module and
$N_{C/Y}=\mathcal O_C$. Put $Z=V(J)\subset C$. We now compute
the actual restriction map
\begin{equation}\label{eq:E14-jet-map}
 H^1(C,\mathcal O_C)\longrightarrow H^1(Z,\mathcal O_Z).
\end{equation}

In the deformation \eqref{eq:E14-elliptic-family} the critical
base near the cusp is the smooth curve
\[
 \Gamma:\quad t=-3x^2,\qquad
 \epsilon=x^2(3+2x).
\]
Its map to the $\epsilon$-disc is finite flat of degree two; its
central fibre is $\operatorname{Spec}\mathbb C[x]/(x^2)$.
The lifted translation fixes its $x$-coordinate. This follows on
the punctured curve because translation fixes the unique node of
each nodal cubic, and then follows at zero by holomorphic
continuation. On the line coordinate it is
$(x,v)\mapsto(x,c(x)v)$ with $c$ a holomorphic nowhere-zero
contraction. Thus its quotient critical locus is exactly the
elliptic family
\[
 \mathcal Z=(\Gamma\times\mathbb C^*)/
       \langle(x,v)\mapsto(x,c(x)v)\rangle.
\]
In particular $H^0(Z,\mathcal O_Z)$ and
$H^1(Z,\mathcal O_Z)$ are both free rank-one modules over
$A=\mathbb C[x]/(x^2)$, and $H^2(Z,\mathcal O_Z)=0$.

Restriction in this proper family is a holomorphic bundle map
$G|_\Gamma\to R^1(\mathcal Z/\Gamma)_*\mathcal O_{\mathcal Z}$.
On $\Gamma\setminus\{0\}$ the calculation of E13 identifies it,
after a nonzero scalar change of target frame, with the row
\[
 \left(1,-\frac{b}{1+x}\right).
\]
The same frame can be chosen at zero. To check this last point,
normalize $C$. Its normalization is a $\mathbb P^1$-bundle over
the elliptic singular curve $E_0$, and its normalization quotient
is $\mathcal O_{E_0}$: the first jet at the cusp transforms with
derivative one under translation on the additive smooth locus.
The normalization exact sequence, together with
$h^1(\mathcal O_C)=2$, $h^2(\mathcal O_C)=1$ and
$h^\bullet(\mathcal O_{\widetilde C})=(1,1,0)$, implies
that $H^1(\mathcal O_C)\to H^1(\mathcal O_{\widetilde C})$
is onto. Restriction from the ruled normalization to its singular
section identifies the latter group with $H^1(\mathcal O_{E_0})$.
Thus restriction to $E_0$ is onto. The common scalar in the row
above cannot have a zero at zero, since both entries would then
vanish and this restriction would be zero. It is a unit. This
proves the asserted holomorphic target-frame comparison.

The equality of these holomorphic maps on the punctured curve
extends over all of $\Gamma$, and in particular over its
nonreduced central fibre. Hence \eqref{eq:E14-jet-map} is,
in the specified bases, the complex-linear map
\[
 \mathbb C^2\longrightarrow A,\qquad
 (a,d)\longmapsto a+d(-i+ix).
\]
Its matrix in $1,x$ has determinant $i$, so it is an isomorphism.
This is the required jet comparison, not an inference from the
rank at nearby reduced fibres.

The exact sequence $0\to J\to\mathcal O_C\to\mathcal O_Z\to0$
now gives
\[
 h^\bullet(C,J)=(0,1,1).
\]
Together with \eqref{eq:E14-log-coho} and the hexagonal quotient,
the tangent exact sequence gives
$h^\bullet(T_Y)=(1,2,2,0)$.

It remains to verify that the actual two-parameter family has
bijective Kodaira--Spencer map. In the above local trivialization
of its total space the map to the base is
\[
 f_epsilon=\frac{y^2-x^3-\epsilon}{1+x}.
\]
The derivative $\partial_\epsilon f$ has image $-1+x$ in
$\mathcal O_C/J=\mathbb C[x]/(x^2)$ along the singular curve.
If the ambient first-order deformation were trivial, every
first-order variation of its map would be a global section of
$f^*T_{\mathbb P^1}$, hence the pullback of a projective vector
field, because $f_*\mathcal O_Y=\mathcal O_{\mathbb P^1}$.
Its local image in this Jacobian quotient is constant. Changing
the local ambient trivialization adds $df(v)$ and hence adds
zero in that quotient. The nonzero coefficient of $x$ therefore
proves that the $\epsilon$-direction is nontrivial and cannot be
cancelled by the scale direction, which does not change this local
map. With $\epsilon=0$, a normalized singular elliptic curve is
$\mathbb C^*/\langle c_0u\rangle$; its modulus has nonzero
$u$-derivative for sufficiently small $u$. A trivial deformation
of the total space preserves its anticanonical pencil from
\eqref{eq:E14-canonical}, its three distinguished special values,
and that elliptic singular curve. Hence the scale direction is
nontrivial too. The two directions span $H^1(T_Y)$.
The inverse-function argument of E13 consequently makes the
Kuranishi ideal zero and proves miniversality. E6 and E1 give the
remaining assertions.
\end{proof}

\begin{statement}{E15}{Proved}{Tangent cohomology and Kuranishi germs of positive combs}
For each positive-comb manifold $X_m$ constructed in
\cite{Viaclovsky2026Combs}, in particular for every $m\geq12$ in
its stated construction, one has
\[
 h^\bullet(X_m,T_{X_m})=(2m+1,\,8m+1,\,6m+1,\,0).
\]
Its Kuranishi germ is smooth and reduced of dimension $8m+1$.
A locally miniversal family is supplied by all sufficiently nearby
admissible period sections on the root curve, with the affine
attachment classes fixed. Its members have algebraic dimension one
and torus fibrations with the same three critical values,
multiplicities $3,4$, and hexagonal fibre type.
\end{statement}
\begin{proof}
Let $\mathfrak B=\mathbb P(3,4)$, with coarse projection
$\rho:\mathfrak B\to B$. Write $p_1,p_2$ for the finite-isotropy
points and $p_0$ for the hexagonal value. On the root curve let
$\mathcal H$ be the extended cohomology bundle, $F$ the Hodge
subbundle, $G$ its quotient, and $V$ its dual:
\[
 0\longrightarrow F\longrightarrow\mathcal H
 \longrightarrow G\longrightarrow0,\qquad V=F^*.
\]
The finite-cover comparison identifies
$G=R^1\bar f_*\mathcal O_{X_m}$ by invariant pushforward of the
ordinary torus Hodge sequence. At the hexagon the two deck
characters of the infinite toric quilt map to the two generators
of $H^1(\mathcal O_W)$ in its normalization complex, establishing
the same comparison there. The bundle $\bar f^*V$ is the vertical
logarithmic tangent bundle: its translations are dual to the
holomorphic one-forms on smooth fibres, and at the hexagon the
dual logarithmic coordinate fields extend in every toric chart.
Cup product identifies $R^2\bar f_*\mathcal O$ with
$D:=\bigwedge^2G$. On finite covers this is the ordinary torus
exterior algebra; at the hexagon it is the exterior algebra of
the two cycles in the normalization complex.

We specify finite characters by the orders of their invariant
generators. At $p_1$ the nontrivial Hodge line has invariant
generator $w$ times an equivariant frame, while the nontrivial
line of each of $V$ and $G$ has generator $w^2$ times its frame.
At $p_2$ the corresponding orders are $3$ for $F$ and $1$ for
$V,G$. Each bundle also has one trivial-character line. These
characters persist because comb smoothing fixes the limiting
periods. The tautological quotient on each tooth is
$\mathcal O(1)^{\oplus2}$, so each tooth adds $2$ to the coarse
degree of $G$. Therefore
\[
 \deg\rho_*G=2m-1,\qquad D=\mathcal O_{\mathfrak B}(24m-1).
\]
The determinant of $\mathcal H$ has degree zero on the root
curve: its determinant local system is trivial, and its cusp
residue has trace zero. Thus
\[
 \det F=\mathcal O(1-24m),\quad
 \det V=\mathcal O(24m-1),\quad
 \deg\rho_*V=2m-1.
\]
Indeed the finite fractional correction for $V$ and $G$ is
$2/3+1/4=11/12$; subtracting it from $(24m-1)/12$ gives
$2m-1$. For $F$ the correction $1/3+3/4=13/12$ gives
$\deg\rho_*F=-2m-1$.

The coarse bundles $V_c=\rho_*V$ and $G_c=\rho_*G$ have
nonnegative splitting. To prove this, evaluate the unchanged
invariant real homology class $\delta$ on $F$. Its resulting
holomorphic section of $V$ never vanishes on a smooth fibre:
a real class annihilating a Hodge plane also annihilates its
conjugate. The same argument holds on the finite good-reduction
covers. At the hexagon, $\delta$ is a nonzero combination of
angular cycles and the logarithmic forms have identity angular
period matrix. Its evaluation is nonzero there too, and gives
a nowhere-zero section of $V_c$. The invariant real cohomology
class $\gamma$ likewise gives a nowhere-zero section of $G_c$;
at the cusp it belongs to $\ker N$, whose projection to the
Hodge quotient is an isomorphism in the transverse graph chart.
At each finite point these nonzero invariant central values give
nonzero values of the coarse sections. A split rank-two bundle
of positive degree on $\mathbb P^1$ with a negative summand
cannot have a nowhere-zero section: every section lies in its
other summand, which has positive degree and hence has zeros.
Since $2m-1>0$, both splittings are nonnegative. In particular
\begin{equation}\label{eq:E15-positive}
 h^0(V_c)=h^0(G_c)=2m+1,\qquad H^1(V_c)=H^1(G_c)=0.
\end{equation}

Set $A=\rho_*(V\otimes G)$ and $L=\mathcal O_{\mathfrak B}(-5)$.
Multiplication of invariant sections gives exact comparisons
\begin{equation}\label{eq:E15-products}
\begin{split}
 0&\longrightarrow V_c\otimes G_c\longrightarrow A
       \longrightarrow\mathbb C_{p_1}\longrightarrow0,\\
 0&\longrightarrow V_c(2m-1)\longrightarrow\rho_*(V\otimes D)
       \longrightarrow\mathbb C_{p_1}\longrightarrow0,\\
 0&\longrightarrow G_c(-1)\longrightarrow\rho_*(L\otimes G)
       \longrightarrow\mathbb C_{p_1}\longrightarrow0.
\end{split}
\end{equation}
Here is the local verification. At $p_1$, multiplying two
nontrivial invariant generators of $V,G$ gives $w^4$, whereas
their tensor character has generator $w$; multiplication is
therefore one factor $t=w^3$. At $p_2$ the product is $w^2$
and the tensor generator also has order two, giving no defect.
Products involving a trivial character have no defect. The
line $D$ has the nontrivial character of $G$ and coarse degree
$2m-1$, proving the second comparison. The invariant-generator
orders of $L$ are $1,1$, and $\rho_*L=\mathcal O_B(-1)$.
Its product with the nontrivial generator of $G$ has the same
single defect at $p_1$, proving the third comparison.

All summands on the left in \eqref{eq:E15-products} have zero
$H^1$, by \eqref{eq:E15-positive}; for the last sequence their
degrees are at least $-1$. Consequently
\[
\begin{array}{c|c|c}
\text{bundle on }B&h^0&h^1\\ \hline
 A&8m+1&0\\
 \rho_*(V\otimes D)&6m&0\\
 \rho_*(L\otimes G)&2m&0
\end{array}
\]
by Riemann--Roch: their degrees are $8m-3$, $6m-2$, $2m-2$
and their ranks $4,2,2$. Also
$\rho_*(L\otimes D)=\mathcal O_B(2m-1)$.
Finite-group invariant pushforward is exact, so these calculations
give the cohomology on the root curve.

The logarithmic tangent sequence is
\[
 0\longrightarrow\bar f^*V\longrightarrow T_{X_m}(-\log W)
 \longrightarrow\bar f^*L\longrightarrow0.
\]
Projection and Leray yield
\[
 h^\bullet(\bar f^*V)=(2m+1,8m+1,6m,0),\qquad
 h^\bullet(\bar f^*L)=(0,2m,2m,0).
\]
Its only possible connecting map between nonzero groups is
$H^0(L\otimes G)\to H^0(V\otimes D)$. This is the map of
sections of the fibrewise cup-product boundary in E5. The cusp
has rank two, so its sheaf map is generically injective. A map
from a locally free sheaf on a smooth curve with zero generic
kernel is injective; hence this map on sections is injective.
It follows that
\[
 h^\bullet(T_{X_m}(-\log W))=(2m+1,8m+1,6m,0).
\]
The branch-normal quotient is $\eta_*K_{dP_6}$, since the toric
three-form trivializes the canonical bundle near the fibre.
Its cohomology is $(0,0,1)$, and its exact sequence proves the
asserted tangent-cohomology tuple.

We construct a family with this full tangent space. All nearby
root-period sections have normal bundle $V\otimes G$, tangent
space $H^0(A)$, and obstruction space $H^1(A)=0$. This section
germ can be realized using an ordinary coarse Grassmannian bundle
and finite analytic conditions. Split the equivariant bundle as
$F\oplus G$ on a finite-cover chart and write a nearby plane as
a graph matrix. Dividing the invariant-generator factors shows
that a holomorphic lift requires one entry to vanish at $p_1$
and two at $p_2$. At $p_1$ the entry is from the nontrivial Hodge
line to the trivial quotient line. At $p_2$ it is that entry
and the entry between the two nontrivial lines. These are linear
conditions on the values of the local graph matrix. Their normal
sheaf is $A$. The Cech gluing obstruction across a small Artin
extension lies in $H^1(A)=0$; each local graph chart and its
stated conditions is smooth. The finite-dimensional Douady
section germ with these conditions is therefore smooth of
dimension $8m+1$. This description does not require a global
manifold cover of the root curve.

These sections produce compact threefolds even though their
limiting periods are allowed to vary. At a finite point the
equivariant lift is a small deformation on the compact
good-reduction cover. Admissibility is open there. The affine
action has the same real formula $x\mapsto Tx+v/n$; it is
holomorphic for the new equivariant periods and stays free
because its real action is unchanged. At the hexagon the
transverse graph chart is open, and its periods are
\[
 Z(q)=\frac{\log q}{2\pi i}I+H(q)
\]
with $H$ holomorphic and uniformly bounded on a smaller parameter
polydisc. The translations with unit factors
$\exp(2\pi iH_{ij})$ preserve the same toric fan. On logarithmic
radii their action is
$r\mapsto r+(I+O(1/\log|q|))\lambda$; it is uniformly proper
and free for small $q$. A fixed finite collection of toric unit
polydiscs gives compact representatives. Thus the filling is
proper with local equation $q=x_1x_2x_3$, and its total space
is smooth over the parameter. Varying $H(0)$ varies gluing units
and preserves these estimates and local equations. On the
compact complement of the critical discs admissibility is open.
The annular gluings extend holomorphically with their fixed
integral classes: the lattice--Lie--torus exact sequence and
Cartan vanishing on an annular product extend the original
translations. This constructs a proper holomorphic family with
all $8m+1$ period directions and the specified central member.

Its Kodaira--Spencer map is injective. A first-order trivialization
would make its map variation a section of
$H^0(X_m,f^*T_B)=H^0(B,T_B)$. The three critical values remain
constant by the equations $w^3,w^4,x_1x_2x_3$, so that base
vector field vanishes at three points and is zero. The
trivialization would then be fibre preserving. On a marked
smooth torus the period derivative is the corresponding
constant Dolbeault tensor in
$\operatorname{Lie}(F)\otimes H^1(F,\mathcal O_F)$.
A nonzero constant tensor cannot be a Dolbeault coboundary,
since derivatives of periodic coefficients integrate to zero.
Thus a period section giving a trivial total-space deformation
vanishes on every regular fibre, and hence identically. Its
Kodaira--Spencer map is an isomorphism onto the already computed
$H^1(T_{X_m})$. The inverse-function argument of E13 then makes
the Kuranishi ideal zero, proving smoothness, reducedness and
miniversality. The topology theorem of the comb construction
identifies these manifolds with $S^6$; E1 and E6 give algebraic
dimension one throughout their small deformation germs.
\end{proof}

\paragraph{Comparison with the constructed parameter space.}
The germ in \cite[Lemma labelled \texttt{paramspace} and
Proposition labelled \texttt{effective}]{Viaclovsky2026Combs}
fixes the values of the coarse period section at all three
critical points. Its normal sheaf is $N(-p_0-p_1-p_2)$, where
$N=(\rho_*F)^*\otimes G_c$. The actual inclusion of that germ
into our root-period germ has the differential induced by
\[
0\longrightarrow N(-p_0-p_1-p_2)\longrightarrow A
\longrightarrow
 \mathbb C_{p_0}^{4}\oplus\mathbb C_{p_1}^{3}
                     \oplus\mathbb C_{p_2}^{2}
\longrightarrow0.
\]
The local graph conditions above prove this exact sequence:
the old sheaf sets all four entries to zero at each point,
whereas $A$ imposes zero, one, and two conditions, respectively.
For the published combs its left term has zero $H^1$, as proved
in the cited smoothing construction. Its tangent dimension is
$8m-8$, and the quotient contributes exactly nine directions.
This explains the relation of the published lower bound to the
dimension $8m+1$ proved here by an explicit morphism of parameter
germs. The source asks for the full Kuranishi germ in
Section~10(1). Searches on 15 November 2026 for its title together
with ``Kuranishi'', ``smooth'', and ``$8m+1$'', and for the
two-parameter and Liu--Rollenske titles with ``tangent cohomology'',
located the supplied construction papers but no earlier complete
computations E13--E15. This records the literature comparison;
we claim no priority.


\begin{statement}{E17}{Refuted}{The embedded-component omission in the divisor-twisting lemma}
CDP98, Proposition~1.1, in its printed generality of arbitrary
nonreduced compact spaces, is false. The exact replacement is as
follows. For an effective Cartier divisor $D$ on a compact complex
space $Z$ and a locally free sheaf $\mathcal V$, the spaces
$H^0(Z,\mathcal V(-kD))$, viewed inside $H^0(Z,\mathcal V)$,
stabilize for $k\gg0$ at
\[
 \{s\in H^0(Z,\mathcal V):\operatorname{supp}(s)\cap|D|
       =\varnothing\}.
\]
In particular they eventually vanish when $Z$ has no embedded
associated components and $D$ meets every irreducible component.
This last case applies to every full Cartier fibre used in E2.
\end{statement}
\begin{proof}
Let $C=\mathbb P^1$, let $p,q$ be distinct points, and form the
square-zero complex space
\[
 |Z|=C,\qquad\mathcal O_Z=\mathcal O_C\oplus\mathbb C_p,
 \qquad(a,u)(b,v)=(ab,av+bu).
\]
Near $p$ its analytic local ring is
$\mathbb C\{z,e\}/(ze,e^2)$, and away from $p$ it is the
ordinary smooth curve. These charts glue through the quotient
by $e$, so this specifies a compact connected complex space.
It has one irreducible component, of positive dimension, and an
embedded associated point at $p$. The divisor $D=q$ is effective
Cartier: its local equation is a coordinate at $q$ and a unit
near $p$. For every integer $k$ there is the exact, split
description as an $\mathcal O_C$-module
\[
 \mathcal O_Z(-kD)=\mathcal O_C(-kq)\oplus\mathbb C_p.
\]
It gives every cohomology group:
\[
 h^0=\max(1-k,0)+1,\qquad h^1=\max(k-1,0),\qquad
 h^j=0\quad(j\geq2).
\]
In particular the nonzero section of the square-zero summand
survives for every $k\geq1$, refuting the printed proposition.

For the replacement, the descending spaces in the statement are
finite dimensional, hence stabilize. A section in their
intersection belongs to $\bigcap_k I_D^k\mathcal V_z=0$ at
every $z\in|D|$, by Krull intersection in the analytic local
ring. Its coherent support is therefore disjoint from $D$.
Conversely a section with support disjoint from $D$ vanishes
near $D$ and is a section of $\mathcal V(-kD)$ for every $k$.
This proves the exact stable-space formula. When there are no
embedded associated components, every nonzero submodule of a
locally free sheaf has support containing an irreducible
component of $Z$. Such a support cannot avoid $D$ under the
stated intersection hypothesis, proving the vanishing.
\end{proof}

\paragraph{Audit of the printed non-holomorphic proof.}
We now give its entire chain of reductions and identify its
dependencies. CDP20, Section~3, first uses constancy of the index
under Picard twists and Riemann--Roch to deduce its parts (b)
and (c) from generic vanishing. Serre duality reduces generic
vanishing to degrees zero and two. Degree zero invokes CDP98,
Corollary~1.3; that corollary uses only Proposition~1.1 on the
smooth manifold, a multiple of an effective divisor lying in
$\operatorname{Pic}^0$, and semicontinuity. E17 proves the
version of Proposition~1.1 needed there, and E2 spells out the
exponential-family semicontinuity argument. Thus this dependency
is valid and involves none of the cotangent or monodromy steps.

For degree two the printed proof resolves the meromorphic pencil
by $\sigma:\widehat Y\to Y$, with a connected map
$h:\widehat Y\to\mathbb P^1$, and writes
\[
 h^*\mathcal O(1)=\sigma^*P\otimes\mathcal O(-E),
\]
where $E$ is the exceptional divisor of the pulled-back base
ideal. The map on second cohomology from $Y$ to $\widehat Y$
is an isomorphism, since a composition of smooth blowups has
$R\sigma_*\mathcal O=\mathcal O_Y$. The divisor sequence
reduces vanishing on $\widehat Y$ to vanishing after $-tE$
and to $H^2(tE,\sigma^*W)=0$. The printed proof asserts
the exceptional higher-image vanishing needed for this step.
For this particular pencil divisor it has the following direct
justification. Pull back the evaluation sequence on $\mathbb P^1$
and tensor by $\sigma^*P^{-t}$ to obtain, for $t\geq1$,
\[
\begin{split}
0\longrightarrow&\bigl(\mathcal O(E)\otimes
              \sigma^*P^{-(t+1)}\bigr)^{\oplus t}
 \longrightarrow(\sigma^*P^{-t})^{\oplus(t+1)}\\
 &\longrightarrow\mathcal O(-tE)\longrightarrow0.
\end{split}
\]
The evaluation sequence remains exact under pullback because it
is a sequence of vector bundles with locally free quotient.
The middle term has zero $R^2\sigma_*$, by the projection
formula and $R\sigma_*\mathcal O=\mathcal O_Y$.
Every coherent sheaf has zero $R^3\sigma_*$ because the
fibres have dimension at most two. Hence
$R^2\sigma_*\mathcal O(-tE)=0$; the divisor sequence now gives
\[
 R^1\sigma_*\mathcal O_{tE}=
 R^2\sigma_*\mathcal O_{tE}=0.
\]
The support of $\sigma_*\mathcal O_{tE}$ has dimension at most
one, so Leray and the projection formula give the required
$H^2(tE,\sigma^*W)=0$. No identification of this pushforward
with a structure sheaf is needed.

This special verification matters. The analogous assertion for
an arbitrary exceptional divisor is false. Blow up a point
$p\in\mathbb P^3$ and then a smooth plane cubic in the exceptional
$\mathbb P^2$. For the second exceptional divisor $S$, projection
to that cubic is a $\mathbb P^1$-bundle; thus
$h^\bullet(S,\mathcal O_S)=(1,1,0)$ and the composite contraction
to $\mathbb P^3$ has
$R^1\sigma_*\mathcal O_S=\mathbb C_p$.
This example, also given in \cite{LRS2019}, concerns arbitrary
exceptional divisors and does not refute the just-proved assertion
for the divisor of a resolved pencil. The graph proof of E2
avoids the distinction by computing its own higher direct images
from its hypersurface equation.

The printed argument then replaces its twist $M$ by
$M\otimes P^{t+k}$ and applies Leray to $h$. It must prove
\[
 R^2h_*\sigma^*(W\otimes M\otimes P^k)=0
\]
and then the vanishing of $H^1$ of the corresponding $R^1$
twisted by $\mathcal O(t)$; the latter is ordinary Serre
vanishing on $\mathbb P^1$. For the former it uses top-degree
base change and fibrewise Serre duality. It invokes CDP98,
Proposition~1.1, first on components meeting $E$, and then
successively on components at increasing graph distance from
them. Coefficients in an auxiliary effective divisor are chosen
in reverse order of this distance, and the power of $E$ is
chosen last. E2 gives this argument with the actual pullback
maps, full multiplicity filtrations, and a finite-support
argument establishing uniformity. These fibres are Cartier in
a smooth threefold and have no embedded components, so E17's
proved version of the twisting lemma applies.

These are the only earlier CDP results invoked in Section~3:
CDP98, Corollary~1.3 and Proposition~1.1. Its opening reference
to CDP98, Theorem~2.1 identifies the proof being rewritten,
rather than using that theorem as a premise. CDP98, Lemma~1.5
is not invoked; its erroneous replacement of fibrewise vanishing
by generic vanishing is precisely what the full component
argument removes. No result from CDP20, Sections~4--7 is used.
In particular the cotangent-filtration and monodromy failures
treated in Parts A and D do not enter E2. The supplied LRS
graph proof has the same dependency on the no-embedded-component
twisting lemma, and supplies its own higher-image calculation.
Its displayed adjunction formula has a sign misprint:
for a Cartier fibre $F\subset Y$ it is
$\omega_F=(K_Y+F)|_F$, not $(K_Y-F)|_F$.
E2 uses the actual dualizing line and includes nonreduced fibres,
so that misprint has no role in the proof here.


\begin{statement}{E18}{Refuted}{The square theta formula for arbitrary sections}
The square-only formula in Lead~9 is false as a local period
identity for arbitrary elliptic sections and linearizations.
The exact formula for the theta expression specified there is
\eqref{eq:E-theta-general}. The assertion for the actual original
period family remains true, by its equation $\wp(\mu)=0$.
\end{statement}
\begin{proof}
Let $\Delta$ be a sufficiently small disc about $i$ in the upper
half-plane and choose the holomorphic logarithm of the nonvanishing
function $\vartheta_1(1/2,\tau)/\eta(\tau)$ which is real at $i$.
Put
\[
 \mu(\tau)=\frac12,\qquad
 v(\tau)=\frac{i}{\pi}
       \log\frac{\vartheta_1(1/2,\tau)}{\eta(\tau)}+i,
 \qquad \beta(\tau)=6v(\tau).
\]
At $\tau=i$, the product expansions give
\[
 \frac{\vartheta_1(1/2,i)}{\eta(i)}
 =2e^{-\pi/6}\prod_{n\geq1}(1+e^{-2\pi n})^2>1.
\]
Thus $\operatorname{Im}v(i)>1$, and after shrinking $\Delta$
one has $\operatorname{Im}v>0$. Define
\[
 \Lambda_\tau=\mathbb Z(1,0)+\mathbb Z(0,1)
 +\mathbb Z(\tau,1/2)+\mathbb Z(3,\beta(\tau)),
 \qquad T_\tau=\mathbb C^2/\Lambda_\tau.
\]
The real determinant of these four periods is nonzero, since
$\operatorname{Im}\tau>0$ and $\operatorname{Im}\beta>0$.
They form a locally uniformly full lattice, so the quotients form
a proper holomorphic family of compact complex two-tori over
$\Delta$. Each fibre has
$h^\bullet(\mathcal O)=(1,2,1)$ and
$h^\bullet(T)=(2,4,2)$, with zero higher groups, by its
translation-invariant Dolbeault complex.

This family also has the required elliptic line-bundle
description. Quotient first by $(1,0),(0,1),(\tau,1/2)$.
Writing $z$ for the first coordinate and $a=\exp(2\pi iz_2)$
for the second gives the complement of zero in the flat line
bundle on $\mathbb C/(\mathbb Z+\tau\mathbb Z)$ with multipliers
$1$ and $-1$. The theta quotient for the point $P=1/2$ identifies
this line bundle with $\mathcal O(P-O)$. The last lattice
generator acts by
$(z,a)\mapsto(z+3,\exp(2\pi i\beta)a)$.
Its elliptic translation is $-6P=0$ in the elliptic curve;
its scalar is a holomorphic contracting linearization. This
constructs the family from an elliptic family, a section and a
linearization explicitly.

There is a specified relation to the symmetric period matrix
$\Omega=\left(\begin{smallmatrix}\tau&\mu\\\mu&v\end{smallmatrix}\right)$.
Adjoining $(\mu,v)$ enlarges $\Lambda_\tau$ to the lattice
generated by $(1,0),(0,1),(\tau,\mu),(\mu,v)$; its index is six,
because $6(\mu,v)=(3,\beta)$ and no smaller multiple belongs
to the original real lattice. Identity on $\mathbb C^2$ induces
a finite holomorphic isogeny of degree six between these families.
Its differential is the identity on the Lie spaces and its
pullback on complex first cohomology is an isomorphism. The
Kodaira--Spencer tensors therefore have the same rank under
these explicitly given comparisons.

The identity $\vartheta_1(1/2+z,\tau)=\vartheta_2(z,\tau)$
and evenness of $\vartheta_2$ give
$\partial_\mu\log\vartheta_1(1/2,\tau)=0$. Since $\mu'=0$,
the proposed square is zero. The actual determinant, by the
proved heat-equation formula, is
\[
 \det\Omega'(i)=-\frac{\wp(1/2;i)}{4\pi^2}
 =-\frac{\vartheta_3(0,i)^4+\vartheta_4(0,i)^4}{12}\ne0.
\]
The last inequality follows from the theta products at the real
nome $e^{-\pi}$, whose factors give strictly positive theta
constants. Thus even the zero sets in the asserted formula
differ in this fully specified family. Extra holomorphic terms
in a different linearization contribute their derivatives to
the determinant; they must be retained, just as the $\wp$ term
must be retained when the section is not a zero of $\wp$.

For completeness, the direct images and total-space cohomology
in this local example are also determined. Write
$g:\mathscr T\to\Delta$ for the displayed proper family.
Translation-invariant forms, base change, and triviality of
holomorphic vector bundles on a disc give
\[
 R^qg_*\mathcal O_{\mathscr T}
   =\mathcal O_\Delta^{\binom2q},\qquad
 R^qg_*T_{\mathscr T/\Delta}
   =\mathcal O_\Delta^{2\binom2q}\quad(0\leq q\leq2),
\]
with all higher images zero. Shrink once more so that the
nonzero determinant stays nonzero. In the direct-image sequence
of $0\to T_{\mathscr T/\Delta}\to T_{\mathscr T}
\to g^*T_\Delta\to0$, the first boundary sends $1$ to the
nowhere-zero tensor $\kappa$, and the next boundary is the
rank-two isomorphism of E5. Thus
\[
 R^qg_*T_{\mathscr T}=
 \begin{cases}
  \mathcal O_\Delta^2&q=0,\\
  \mathcal O_\Delta^3&q=1,\\
  \mathcal O_\Delta&q=2,\\
  0&q>2.
 \end{cases}
\]
Since $\Delta$ is Stein, every total-space group
$H^q(\mathscr T,\mathcal F)$ for these three sheaves is exactly
the space of sections on $\Delta$ of the corresponding displayed
direct image. This specifies every cohomological degree without
assigning finite dimensions to the holomorphic-function spaces
on a disc.
\end{proof}


\begin{statement}{E19}{Proved}{Homology of all original affine twists}
For the original multiplicities $(3,4)$ and every admissible affine
datum, put $p=12l_0-4r_1-3r_2$ as in R5 and $d=|p|$.
Then $p\ne0$ and
\[
 \pi_1(X)=\mathbb Z/d\mathbb Z,
 \qquad
 H_j(X,\mathbb Z)=
 \begin{cases}
  \mathbb Z&j=0,6,\\
  \mathbb Z/d\mathbb Z&j=1,2,3,4,\\
  0&j=5.
 \end{cases}
\]
In particular every $|p|>1$ member is a rational homology
six-sphere, has algebraic dimension one, and has the deformation
and tangent-cohomology conclusions of E11.
\end{statement}
\begin{proof}
We check and extend the gluing calculation of
\cite[Propositions labelled \texttt{mvred}, \texttt{mumford},
and \texttt{matrix}]{Viaclovsky2026Log}. Its section on homology
assumes that the gluing determinant is $\pm1$. We show explicitly
which argument works for every nonzero determinant.

First specify the marking and the collar change. In the notation
of the conventions, the change from the retained lattice to the
ordered lattice $(\delta,a,b,d_0)$ of that paper is
\[
 R=\begin{pmatrix}
 -2&1&0&1\\-6&0&-1&0\\0&-1&0&0\\-1&0&0&0
 \end{pmatrix},\qquad \det R=-1.
\]
Its intertwining identities are $A_0R=RA_1^{\rm retained}$ and
$A_1R=RA_2^{\rm retained}$. With $\psi=d_0^\vee$ one has
$\psi R=-\gamma$. The inverse affine-generator comparison
therefore sends an invariant translation $w$ to $-Rw$.
In particular, for $u_0=d_0+2a+2b$ and
$u_1=d_0+3a+3b$,
\[
 -R\varepsilon=u_0,\qquad
 -R\varepsilon'=u_1-\delta,\qquad
 -R\widehat\delta=-\delta.
\]
Thus the coefficients of $u_0,u_1$ before changing the cusp
reference section are exactly $r_1,r_2$.

A continuous reference section on the punctured base can be
chosen to agree with the section of the Mumford filling on its
annular collar. Here is the extension comparison. The punctured
base retracts onto a graph, so its torus torsor has obstruction
in $H^2(U,\Lambda)=0$. A difference of sections on the cusp
annulus has a winding class in
$H^1(A_\infty,\Lambda)=\Lambda/(A_\infty-I)\Lambda$.
Restriction from graph cocycles is onto: the cocycle
$(v_0,v_1)$ on its two edges restricts to
\[
 -A_1^{-1}v_1-A_\infty v_0
       \pmod{(A_\infty-I)\Lambda}.
\]
Taking $v_0=0$ proves surjectivity. Smooth sections of the Lie
bundle form a fine sheaf, so every such winding class is realized
by a torus section. The remaining difference has a Lie lift,
which extends by a cutoff on the collar. This constructs the
required reference section and identifies the underlying real
torus bundle with the linear flat bundle using that section as
zero.

In the retained peripheral presentation write $c$ for the
primitive common coinvariant and $x,y$ for the finite meridians.
The relations are $x^3=c^{r_1}$, $y^4=c^{r_2}$ and
$xy=c^{l_0}$, with $c$ central. Replacing
$y$ by $c^{-l_0}y$ is the exact change to zero cusp winding.
It gives finite coefficients
\[
 n_0=r_1,\qquad n_1=r_2-4l_0,
 \qquad \Delta=4n_0+3n_1=-p.             \tag{E.19.1}
\]
The remaining winding component and the other affine translations
only change the coefficients $\rho_i$ of $\delta$ in
$v_i=\rho_i\delta+n_i u_i$. They will be retained in the
gluing matrix. Admissibility gives
$\gcd(n_0,3)=\gcd(n_1,4)=1$, so $\Delta$ is odd and is not
divisible by three. In particular it is nonzero.

The actual van Kampen maps kill $\delta,a$ and the cusp meridian.
The relation sending $a$ to $-b$ then kills $b$. The remaining
generators $c,x$ commute and have relations
\[
 -n_0c+3x=0,\qquad -n_1c-4x=0.
\]
The relation matrix has determinant $\Delta$ and first Smith
divisor one, because its entries include $3,4$. This proves the
asserted fundamental group and first homology.

Let $V$ be the complement of the three filling interiors and
$M_i$ their boundaries. Each $H_1(M_i)$ is free of rank three;
$H_1(V)$ is free of rank three; and each reduced filling fibre
$F_i$ has free $H_1$ of rank two. At a finite fibre the last
claim also follows directly from
\[
 H_1(F_i)=\mathbb Z\langle a,d_0,g_i\rangle/
  \mathbb Z(s_i\rho_i a-n_i d_0+m_i g_i),
 \quad(s_0,s_1)=(3,2),\quad(m_0,m_1)=(3,4),
\]
whose relation is primitive by the admissibility just checked.
In the Mayer--Vietoris sequence the degree-one map is therefore
between free rank-nine groups and has finite cokernel
$H_1(X)$. It is injective. This is the only place the source
used $H_1(X)=0$ to eliminate that term; finite first homology
suffices.

The map $H_2(M_i)\to H_2(F_i)$ is surjective. The subsequent
pair map takes a cycle to its intersection with $F_i$, and this
intersection is zero because $m_iF_i$ is principal. The pair
sequence and Poincare--Lefschetz duality thus give
\[
 H_2(X)=\operatorname{coker}\left(
  \bigoplus_i H^3(F_i)\xrightarrow{\sum\tau_i}H_2(V)
 \right),
\]
where $\tau_i$ sends a loop on $F_i$ to its normal-circle torus.
All these identifications are integral.

For clarity the basis calculation in this cokernel is also
checked. The graph spectral sequence gives
\[
0\longrightarrow(\textstyle\bigwedge^2\Lambda)_\Gamma
 \longrightarrow H_2(V)
 \longrightarrow\ker[(A_0-I),(A_1-I)]\longrightarrow0.
\]
The first term is generated freely by $s=[a\wedge b]$:
the exterior-square relations kill
$\delta\wedge a,\delta\wedge b,a\wedge d_0,b\wedge d_0$
and give $\delta\wedge d_0=-6a\wedge b$; the invariant
form $\xi=a^\vee\wedge b^\vee+6d_0^\vee\wedge\delta^\vee$
evaluates to one on the remaining generator. The kernel on the
right is free of rank five. Its basis consists of the two cusp
vectors $-(A_1\lambda,\lambda)$ for $\lambda=\delta,a$ and
the three vectors $(\delta,0),(u_0,0),(0,u_1)$.
The minor in rows $1,3,4,5,8$ of these five columns is $-1$,
so they form a primitive rank-five subgroup and hence the full
integral kernel.
The two cusp vectors are the boundaries of the section disc
swept by its two extending circle actions. The normalization
sequence gives $H^3(F_\infty)=\mathbb Z^2$, so they generate
its entire image. Removing them leaves the free basis
$(s,t,t_0,t_1)$ of a rank-four quotient.

Choose $m_i\alpha_i+n_i\beta_i=1$ and
$z_i=\alpha_i d_0+\beta_i g_i$. The domain basis is
$(a,z_0,a,z_1)$. A normal-circle torus over a lattice loop
$\lambda$ has fibre component
$\mathcal N_i\lambda$, with
$\mathcal N_i=\sum_{j=0}^{m_i-1}A_i^j$, and its integral
against $\xi$ is $-\xi(\lambda,v_i)$. For $g_i$ these
quantities are $v_i$ and zero, since $m_i g_i=v_i$ and
$\xi(v_i,v_i)=0$. Substitution of
\[
 \mathcal N_0a=-\delta,\quad\mathcal N_1a=-2\delta,
 \quad\mathcal N_0d_0=3u_0,\quad\mathcal N_1d_0=4u_1
\]
and $\xi(a,u_0)=2$, $\xi(a,u_1)=3$ gives the actual matrix
\[
 B=\begin{pmatrix}
 -2n_0&-6\rho_0\alpha_0&-3n_1&-6\rho_1\alpha_1\\
 -1&\rho_0\beta_0&2&-\rho_1\beta_1\\
 0&1&0&0\\0&0&0&1
 \end{pmatrix}.                              \tag{E.19.2}
\]
The last two rows remove the second and fourth columns by
unimodular operations. The remaining two-by-two block has
determinant $-\Delta$ and contains $-1$. Consequently the Smith
form of $B$ is $\operatorname{diag}(1,1,1,|\Delta|)$.
This proves $H_2(X)=\mathbb Z/d\mathbb Z$.

The Euler number is two: smooth and multiple torus quotients
contribute zero, and the single hexagonal fibre contributes two.
Rational Poincare duality now gives $b_3=0$, since $b_1=b_2=0$.
Integral duality and the universal coefficient sequence give
$H_5=0$, $H_4=\operatorname{Ext}(H_1,\mathbb Z)$ and
$H_3=\operatorname{Ext}(H_2,\mathbb Z)$, proving the remaining
groups. E1 applies to its torus fibration because $b_2=0$;
E11 already computes its tangent cohomology and full miniversal
line without using this topology.
\end{proof}

